\documentclass[a4paper,11pt,reqno]{amsart}

\usepackage[margin=1in]{geometry}
\usepackage[T1]{fontenc}
\usepackage{amsmath,amssymb,mathtools}
\usepackage{booktabs}
\usepackage{array}
\usepackage{longtable}
\IfFileExists{newtxtext.sty}{%
  \usepackage{newtxtext,newtxmath}%
}{%
  \usepackage{lmodern}%
}
\usepackage{microtype}
\usepackage[numbers,sort&compress]{natbib}
\usepackage{enumitem}
\usepackage{needspace}
\usepackage{etoolbox}
\usepackage{xurl}
\usepackage[hidelinks,pdfencoding=auto,psdextra]{hyperref}
\usepackage{bookmark}

\setcounter{secnumdepth}{2}
\setcounter{tocdepth}{1}
\setlength{\parindent}{1.5em}
\setlength{\parskip}{0pt}
\setlength{\emergencystretch}{3em}
\setlist[itemize]{leftmargin=2em,itemsep=0.12em,topsep=0.35em}
\setlist[enumerate]{leftmargin=2.2em,itemsep=0.12em,topsep=0.35em}
\allowdisplaybreaks[2]
\raggedbottom
\Urlmuskip=0mu plus 1mu
\frenchspacing
\tolerance=1800
\pretolerance=500
\ifdefined\pdfgentounicode
  \pdfgentounicode=1
\fi

\theoremstyle{plain}
\newtheorem{theorem}{Theorem}[section]
\newtheorem{proposition}[theorem]{Proposition}
\newtheorem{lemma}[theorem]{Lemma}
\newtheorem{corollary}[theorem]{Corollary}
\newtheorem{conjecture}[theorem]{Conjecture}
\newtheorem*{wowconjecture}{Conjecture (WOW-284)}
\theoremstyle{definition}
\newtheorem{definition}[theorem]{Definition}
\theoremstyle{remark}
\newtheorem{remark}[theorem]{Remark}

\BeforeBeginEnvironment{theorem}{\Needspace{6\baselineskip}}
\BeforeBeginEnvironment{proposition}{\Needspace{6\baselineskip}}
\BeforeBeginEnvironment{lemma}{\Needspace{6\baselineskip}}
\BeforeBeginEnvironment{corollary}{\Needspace{5\baselineskip}}
\BeforeBeginEnvironment{definition}{\Needspace{5\baselineskip}}
\BeforeBeginEnvironment{remark}{\Needspace{4\baselineskip}}
\BeforeBeginEnvironment{proof}{\Needspace{4\baselineskip}}

\newcommand{\F}{\mathbb F_5}
\newcommand{\one}{\mathbf 1}
\newcommand{\Spec}{\operatorname{Spec}}
\newcommand{\tr}{\operatorname{tr}}
\newcommand{\diam}{\operatorname{diam}}
\newcommand{\resultbox}[1]{\boxed{#1}}
\newcommand{\RepoTag}{v2.3.0}
\newcommand{\codefile}[1]{%
  \href{https://github.com/SamPetkov/wow284/blob/\RepoTag/#1}{\path{#1}}}
\newcommand{\datafile}[1]{%
  \href{https://github.com/SamPetkov/wow284/blob/\RepoTag/#1}{\path{#1}}}

\hypersetup{
  pdftitle={Counterexamples, Spectral Obstructions, and Deletion Stability for WOW-284},
  pdfauthor={Samuil Petkov},
  pdfsubject={Counterexamples and structural spectral theory for the minimum-dual-degree conjecture WOW-284},
  pdfkeywords={distance spectrum, dual degree, Moore graph},
  pdfcreator={LaTeX with amsart},
  pdfproducer={pdfTeX},
  pdfdisplaydoctitle=true
}

\title[Counterexamples and obstructions for WOW-284]
{Counterexamples, Spectral Obstructions, and Deletion Stability for WOW-284}
\author{Samuil Petkov}
\address{Department of Physics, \'Ecole normale sup\'erieure,
Universit\'e PSL, Paris, France}
\email{samuil.petkov@phys.ens.psl.eu}
\date{}
\subjclass[2020]{Primary 05C50; Secondary 05C12, 05C35, 05E30}
\keywords{distance spectrum, dual degree, Moore graph}

\begin{document}

\begin{abstract}
WOW-284 asserts that the minimum dual degree of every connected graph of
order at least three and girth at least five does not exceed the negative of
its least distance eigenvalue.  We refute it with exact counterexamples of
orders \(38,39,40,42\), and \(50\), and develop a structural theory of the
failure.  For a connected \(k\)-regular graph of girth at least five and
diameter three, we prove
\[
 \delta^*(G)+\lambda_{\min}(D(G))
 =2k-2-\max_{\theta\ne k}(\theta+1)^2.
\]
Consequently every regular strict counterexample has degree at least six and
diameter at most four, while diameter four forces degree at least ten.  We
solve the associated one-variable nonbacktracking linear program exactly,
including optimizer rigidity.  The optimizer yields a positive-semidefinite
slack matrix whose integral excess gives the stronger universal bound
\[
 n\le\left\lfloor
 \frac{3(k+2)^2(k^2+3)}{18k+41}
 \right\rfloor;
\]
this follows from a three-to-one quantization theorem for the integral excess.
The slack matrix's principal minors also recover local cycle constraints.  In
particular, degree-six counterexamples have order at most \(50\), and at the
remaining boundary the associated signed complement is necessarily
disconnected.  We determine the distance spectra of one- and two-vertex
punctures of Moore graphs and establish a uniform deletion-stability bound:
every deletion of at most five vertices from the Hoffman--Singleton graph
remains a strict counterexample, whereas an explicit six-vertex deletion does
not.  All
theorem-level computations use exact arithmetic.  Lean 4.31 kernel-checks the
full \(50\)-vertex proof,
finite spectral certificates at orders \(38,39,40,42\), and the analytic LP
optimum and rigidity for every integer \(k\ge4\).
\end{abstract}

\maketitle

\section{Introduction}\label{sec:introduction}

Let \(G\) be a finite simple connected graph.  Write \(d(v)\) for the degree
of \(v\), \(N(v)\) for its open neighbourhood, and, for
\(S\subseteq V(G)\), let \(G-S\) denote the subgraph induced by
\(V(G)\setminus S\).  Define
\[
 d^*(v)=\frac1{d(v)}\sum_{u\in N(v)}d(u),
 \qquad
 \delta^*(G)=\min_{v\in V(G)}d^*(v).
\]
The distance matrix is \(D(G)=(d_G(u,v))_{u,v\in V(G)}\).  Its eigenvalues are
ordered as
\[
 \partial_1(G)\ge\cdots\ge\partial_n(G),
 \qquad \partial_n(G)=\lambda_{\min}(D(G)).
\]
Aouchiche and Hansen record the following Graffiti conjecture as
Conjecture~7.16 and attribute it to Fajtlowicz's 1998
\emph{Written on the Wall} report
\cite{Fajtlowicz1998} (see also \cite[Conjecture~7.16]{AouchicheHansen2014}).

\begin{wowconjecture}\label{conj:wow284}
If \(G\) has order at least three and girth at least five, then
\[
 \delta^*(G)\le-\lambda_{\min}(D(G)).
\]
\end{wowconjecture}

For graphs in the domain of the conjecture, put
\[
 \Phi(G)=\delta^*(G)+\lambda_{\min}(D(G)).
\]
Thus \(G\) is a strict counterexample precisely when \(\Phi(G)>0\).

The initial disproof is short.  A degree-\(k\) Moore graph of diameter two has
\[
 A^2=(k-1)I-A+J,
 \qquad
 D=2J-2I-A,
\]
and hence
\[
 \delta^*(G)=k,
 \qquad
 \lambda_{\min}(D)=-\frac{3+\sqrt{4k-3}}2.
\]
The conjecture holds on these graphs exactly for \(k\le3\), with equality at
\(k=3\), and fails for every realizable \(k>3\).  The degree-seven
Hoffman--Singleton graph therefore gives a gap of three.

The purpose of this paper is not merely to list descendants of this graph.  It
addresses three structural questions.

\begin{enumerate}
\item Which spectral mechanism governs regular counterexamples?
\item How restrictive are the degree, diameter, and order conditions?
\item How stable is the counterexample property under deletion?
\end{enumerate}

Our main conclusions are as follows.

\begin{itemize}
\item In diameter three, WOW-284 is equivalent to a two-sided adjacency
spectrum window centred at \(-1\); see Theorem~\ref{thm:diameter-three-score}.
\item Every regular strict counterexample has degree at least six and diameter
at most four.  A diameter-four example, if one exists, has degree at least ten;
see Theorems~\ref{thm:regular-degree-six},
\ref{thm:endpoint-diameter}, and~\ref{thm:diameter-four}.
\item The standard one-variable nonbacktracking linear-programming hierarchy
has exact ceiling
\[
 B_k=\frac{(k+2)(k^2+3)}6,
\]
with a unique optimizer up to positive scaling; see
Theorem~\ref{thm:lp-ceiling}.
\item The optimizer defines one positive-semidefinite slack matrix.  Its trace
 is the LP defect.  Its integral excess satisfies a three-to-one quantization
 inequality, giving
\[
 n\le\left\lfloor
 \frac{3(k+2)^2(k^2+3)}{18k+41}
 \right\rfloor;
\]
 see Theorems~\ref{thm:integral-slack} and~\ref{thm:three-to-one}.  Its
 equality case is arithmetically rigid:
 \[
  (k,n,r)=(103,185220,61740);
 \]
 see Corollary~\ref{cor:three-to-one-equality}.  The same slack matrix's
\(2\times2\) minors yield a general cycle-divisibility sieve and the local
order-\(51\) contradiction
\[
 5N_5=153\cdot11=1683,
\]
while the global theorem gives the degree-\(7,8,9\) order windows
\(74,108,150\); see Corollary~\ref{cor:low-degree-windows}.
\item At the unresolved degree-six order-\(50\) boundary, the integral
signed-complement Gram matrix has rank at least \(30\), and its underlying
signed graph is disconnected; see Proposition~\ref{prop:order50-minus-two}
and Theorem~\ref{thm:order50-disconnected}.
\item One-vertex, adjacent-pair, and nonadjacent-pair deletions of Moore graphs
admit exact invariant-subspace decompositions for their recomputed distance
matrices; see Section~\ref{sec:punctures}.
\item Every deletion of at most five vertices from the Hoffman--Singleton graph
remains a strict counterexample, and this universal radius is sharp; see
Theorem~\ref{thm:hs-radius}.
\end{itemize}

The proofs form two complementary hierarchies.  In the obstruction direction,
the distance-polynomial identity converts WOW-284 into a shifted adjacency
window; scalar trace moments give the one-point LP ceiling; integrality of the
optimal slack matrix strengthens the order bound; graph realizability and
 small Gram minors quantize the remaining excess into the three-to-one bound;
edge-local \(2\times2\) minors convert the same certificate into cycle counts;
and \(3\times3\) minors impose the two-path constraints at order fifty, where
a signed-root representation forces a nontrivial component decomposition.
In the stability
direction, deleting vertices from a Moore graph produces an incidence-Gram
correction to the distance matrix and a configuration-sensitive perturbation
bound.
Invariant-subspace decompositions then give exact puncture spectra, while
orbitwise positive-definiteness certificates determine the sharp
Hoffman--Singleton deletion radius.

The distance-polynomial viewpoint is established in the theory of minimal
cages, distance-polynomial graphs, and spectral excess
\cite{HowladerPanigrahi2022,Fiol2016}.  Nonbacktracking linear-programming
bounds are due to Nozaki and related spectral-Moore work
\cite{Nozaki2015,CioabaEtAl2016}.  Our contribution is the specialization to
the two-sided WOW window, the exact method ceiling, the integral
optimal-slack hierarchy, the edge-local cycle certificate, and the deletion
theory developed below.  Completed in its original counterexample form on
19 July 2026, this work is, to the author's knowledge, the first proof that WOW-284 is false
and the first Lean formalization of such a refutation.  We
do not claim that order \(38\) is minimum or that the constructions classify
all counterexamples.
The literature boundaries are returned to in Section~\ref{sec:scope}.

The analytic arguments are proved in the text.  Precisely specified finite
classifications and matrix certificates are treated as computer-assisted proof
components.  Their exact reproducibility materials are archived with the
accompanying release.

\section{Local growth, Moore graphs, and explicit counterexamples}
\label{sec:examples}

\subsection{Dual degree as radius-two growth}

For \(v\in V(G)\), write \(\Gamma_i(v)\) for the distance-\(i\) sphere and
\(B_2(v)=\Gamma_0(v)\cup\Gamma_1(v)\cup\Gamma_2(v)\).

\begin{proposition}[Backelin's second-degree identity]\label{prop:radius-two}
If \(G\) contains no triangle and no \(4\)-cycle, then
\[
 |B_2(v)|=1+\sum_{u\in N(v)}d(u),
 \qquad
 d^*(v)=\frac{|B_2(v)|-1}{d(v)}.
\]
\end{proposition}

\begin{proof}
For distinct neighbours \(u,w\) of \(v\), the sets
\(N(u)\setminus\{v\}\) and \(N(w)\setminus\{v\}\) are disjoint; an
intersection would form a \(4\)-cycle, and a member in \(N(v)\) would form a
triangle.  These sets partition \(\Gamma_2(v)\), so
\[
 |\Gamma_2(v)|=\sum_{u\in N(v)}(d(u)-1).
\]
Adding the centre and first sphere proves the formula.  This is the normalized
form of Backelin's second-degree lemma \cite[Lemma~2.1]{Backelin2015}.
\end{proof}

\subsection{The Moore threshold}

\begin{theorem}\label{thm:moore-threshold}
Let \(M\) be a degree-\(k\) Moore graph of diameter two, \(k\ge2\).  Then
\[
 |V(M)|=k^2+1,
 \qquad g(M)=5,
 \qquad \delta^*(M)=k,
\]
\[
 \lambda_{\min}(D(M))=-\frac{3+\sqrt{4k-3}}2,
\]
and
\[
 \Phi(M)=k-\frac{3+\sqrt{4k-3}}2.
\]
Thus \(M\) satisfies WOW-284 exactly for \(k\le3\), with equality exactly at
\(k=3\).
\end{theorem}

\begin{proof}
The Moore bound is attained, so adjacent vertices have no common neighbour and
nonadjacent vertices have exactly one.  Given an edge \(uv\), choose
\(x\in N(u)\setminus\{v\}\) and \(y\in N(v)\setminus\{u\}\).  The vertices
\(x,y\) are nonadjacent, since an edge would create a four-cycle, and their
unique common neighbour completes a five-cycle through \(uv\).  Therefore
\[
 A^2=(k-1)I-A+J.
\]
On \(\one^\perp\), the nonprincipal adjacency eigenvalues are the roots
\[
 r,s=\frac{-1\pm\sqrt{4k-3}}2.
\]
Every nonedge has distance two, hence \(D=2J-2I-A\).  The least distance
eigenvalue is \(-2-r=-(3+\sqrt{4k-3})/2\).  Regularity gives
\(\delta^*=k\), and
\[
 (2k-3)^2-(4k-3)=4(k-1)(k-3)
\]
gives the threshold.  The exact scalar and finite checks are independently
repeated by \codefile{scripts/verify_regular_score_calculus.py}.
\end{proof}

\subsection{A coordinate Hoffman--Singleton certificate}

All subscripts below lie in \(\F=\mathbb Z/5\mathbb Z\).  Let
\[
 V(M)=\{P_{i,j}:i,j\in\F\}\mathbin{\dot\cup}
      \{Q_{k,\ell}:k,\ell\in\F\},
\]
with edges
\begin{align*}
 P_{i,j}&\sim P_{i,j\pm1},\\
 Q_{k,\ell}&\sim Q_{k,\ell\pm2},\\
 P_{i,j}&\sim Q_{k,ik+j}.
\end{align*}
This is Hafner's affine-coordinate form of the Hoffman--Singleton graph after a
minor reindexing \cite{Hafner2003}.

\begin{proposition}\label{prop:hs-coordinate}
The coordinate construction is a simple connected \(7\)-regular graph on
\(50\) vertices.  Adjacent pairs have no common neighbour and nonadjacent pairs
have exactly one.  Consequently it has girth five, diameter two, and
\[
 \Spec D(M)=\{91^{(1)},1^{(21)},(-4)^{(28)}\}.
\]
In particular, \(\delta^*(M)=7\) and \(\Phi(M)=3\).
\end{proposition}

\begin{proof}
The neighbourhoods are
\begin{align*}
N(P_{i,j})&=\{P_{i,j-1},P_{i,j+1}\}
 \cup\{Q_{k,ik+j}:k\in\F\},\\
N(Q_{k,\ell})&=\{Q_{k,\ell-2},Q_{k,\ell+2}\}
 \cup\{P_{i,\ell-ik}:i\in\F\}.
\end{align*}
They have seven distinct entries.  For two vertices of the same type, common
neighbours are determined by one nonzero linear equation over \(\F\); for a
cross pair \(P_{i,j},Q_{k,\ell}\), put
\(r=\ell-(ik+j)\).  The pair is adjacent for \(r=0\), has one common
\(P\)-neighbour for \(r\in\{\pm1\}\), and one common \(Q\)-neighbour for
\(r\in\{\pm2\}\).  The five residues are exhausted.  The claimed geometry
and spectrum now follow from Theorem~\ref{thm:moore-threshold}.  The exhaustive
pair certificate, integer BFS distances, characteristic polynomial, and exact
positive-definiteness check are in \codefile{scripts/verify_exact.py}.
\end{proof}

\subsection{Smaller exact counterexamples}

Let
\[
 \mathcal P=\{P_{0,j},Q_{0,j}:j\in\F\}.
\]
The induced graph \(M[\mathcal P]\) is a Petersen graph.  Put
\[
 R=M-\mathcal P,
 \quad H_{39}=R-P_{1,0},
 \quad H_{38}=R-\{P_{1,0},P_{1,1}\},
\]
and let \(X_{42}\) be the second subconstituent of \(P_{0,0}\), namely the
graph induced by the vertices at distance two from it.

\begin{theorem}\label{thm:explicit-examples}
The following are strict counterexamples.
\[
\begin{array}{c@{\quad}c@{\quad}c@{\quad}c}
\toprule
G&|V(G)|&\delta^*(G)&\lambda_{\min}(D(G))\\
\midrule
H_{38}&38&17/3&-3-\sqrt7\\
H_{39}&39&35/6&>-35/6\\
R&40&6&-5\\
X_{42}&42&6&-5\\
M&50&7&-4\\
\bottomrule
\end{array}
\]
The entry for \(H_{39}\) records an exact strict lower bound obtained from
positive definiteness of \(6D+35I\); it is not a decimal approximation to the
least eigenvalue.  Moreover, all \(40\) labelled singleton deletions of \(R\),
and all \(120\)
labelled deletions of the endpoints of an edge of \(R\), are strict
counterexamples with the same respective characteristic polynomials.
\end{theorem}

\begin{proof}
The graph \(R\) is the classical \((6,5)\)-cage of O'Keefe and Wong
\cite{OKeefeWong1979,Wong1979}; its realization as a Petersen deletion of
the Hoffman--Singleton graph also appears in \cite{KlinMuzychukZivAv2009}.
The Moore block identity gives
\[
 \Spec A(R)=\{6^{(1)},2^{(18)},1^{(4)},(-2)^{(5)},(-3)^{(12)}\},
\]
and Theorem~\ref{thm:diameter-three-score} below maps this to
\[
 \Spec D(R)=\{75^{(1)},3^{(5)},0^{(16)},(-5)^{(18)}\}.
\]
The second-subconstituent calculation gives
\[
 \Spec D(X_{42})=\{81^{(1)},4^{(6)},0^{(14)},(-5)^{(21)}\}.
\]
The classical second-subconstituent identification and adjacency spectrum are
recorded in \cite{vanDamHaemers2003}.
For \(H_{38}\), a direct degree count gives \(\delta^*=17/3\), while the
factor \(x^2+6x+2\), together with an exact Sturm isolation, gives the least
root \(-3-\sqrt7\).  For \(H_{39}\), the exact matrix
\(6D+35I\) is positive definite.  All graph, girth, distance, dual-degree,
Sturm, and rational \(LDL^{\mathsf T}\) certificates are checked by
\codefile{scripts/verify_extended.py}.  More explicitly, for every
\(v\in V(R)\),
\begin{align*}
\det(xI-D(R-v))=P_{39}(x):={}&x^9(x+5)^{12}(x^2+6x+3)\\
&\cdot(x^3+3x^2-15x-7)^2(x^3+3x^2-15x-3)^2\\
&\cdot(x^4-78x^3+303x^2-70x-450),
\end{align*}
whereas for every \(uv\in E(R)\),
\begin{align*}
\det(xI-D(R-\{u,v\}))=P_{38}(x):={}&x^4(x-2)(x+5)^8(x^2+6x+2)^2\\
&\cdot(x^3+3x^2-15x-3)\\
&\cdot(x^4+5x^3-7x^2-23x-6)\\
&\cdot(x^5+9x^4+7x^3-77x^2-54x-4)\\
&\cdot(x^9-67x^8-404x^7+1772x^6+7205x^5\\
&\hspace{3em}-18489x^4-17018x^3+20288x^2+16824x+1680).
\end{align*}
The labelled deletion families are exhausted by
\codefile{scripts/verify_descendant_families.py}.  No transitivity or
numerical root ordering is used.  Fixed graph6, adjacency-list, edge-list, and
distance-matrix records are provided in \path{data/graphs/}.
\end{proof}

\begin{theorem}[Moore second subconstituents]\label{thm:second-subconstituent}
Let \(M\) be a degree-\(K\) Moore graph of diameter two, \(K\ge3\), and let
\(X\) be the graph induced by \(\Gamma_2(v)\) for a fixed vertex \(v\).  Then
\(X\) has order \(K(K-1)\), degree \(K-1\), girth at least five, diameter
three, and
\[
 \lambda_{\min}(D(X))=-\frac{5+\sqrt{4K-3}}2.
\]
Consequently
\[
 \Phi(X)=K-1-\frac{5+\sqrt{4K-3}}2,
\]
which is positive exactly for integers \(K\ge6\).
\end{theorem}

\begin{proof}
Relative to \(\{v\}\sqcup N(v)\sqcup\Gamma_2(v)\), write the adjacency
matrix with incidence block \(C\) and induced block \(B\).  The Moore identity
gives
\[
 CC^{\mathsf T}=(K-1)I,
 \qquad CB=J-C,
\]
and on \(\ker C\),
\[
 B^2+B-(K-1)I=0.
\]
Since \(CC^{\mathsf T}=(K-1)I\), the map \(C^{\mathsf T}\) is injective,
its image is orthogonal to \(\ker C\), and
\[
 \mathbb R^{\Gamma_2(v)}
 =\langle\one\rangle\perp C^{\mathsf T}(\one^\perp)\perp\ker C.
\]
Thus \(B\) has principal eigenvalue \(K-1\), eigenvalue \(-1\) on
\(C^{\mathsf T}(\one^\perp)\), and the two Moore roots
\((-1\pm\sqrt{4K-3})/2\) on \(\ker C\).  It remains to identify the
diameter of \(X\).  If the unique common neighbour in \(M\) of two
nonadjacent vertices \(x,y\in X\) lies in \(X\), their distance in \(X\)
is two.  Otherwise it is their common parent in \(N(v)\).  Choose
\(b\in N_X(x)\).  Then \(b\not\sim y\), and the unique common neighbour
\(c\) of \(b,y\) belongs to \(X\): it cannot be \(v\), and it cannot
lie in \(N(v)\), since \(b\) and \(y\) have different parents there.
Thus \(x-b-c-y\) is a path in \(X\).  Pairs with a common parent have no
length-two path in \(X\), so \(X\) has diameter three and
Theorem~\ref{thm:diameter-three-score} applies.  Among the nonprincipal
adjacency eigenvalues, \((-1+\sqrt{4K-3})/2\) maximizes
\(|\theta+1|\); substitution gives
\[
 \lambda_{\min}(D(X))=-\frac{5+\sqrt{4K-3}}2.
\]
The threshold reduces to
\[
 (2K-7)^2-(4K-3)=4(K^2-8K+13)>0
\]
for \(K\ge6\).  The finite \(K=7\) instance is checked by
\codefile{scripts/verify_wow284_38_40_42.py}.
\end{proof}

\section{The regular score calculus}\label{sec:score-calculus}

\begin{theorem}[Diameter-three score formula]\label{thm:diameter-three-score}
Let \(G\) be connected, \(k\)-regular, of girth at least five and diameter
three, with adjacency matrix \(A\) and order \(n\).  Then
\[
 D=3J+(k-3)I-2A-A^2,
\]
\[
 D+kI=3J+(2k-2)I-(A+I)^2.
\]
The principal distance eigenvalue is \(3n-k^2-k-3\), and a nonprincipal
adjacency eigenvalue \(\theta\) gives the distance eigenvalue
\[
 \mu(\theta)=k-2-(\theta+1)^2.
\]
Consequently
\[
 \resultbox{
 \Phi(G)=2k-2-\max_{\theta\ne k}(\theta+1)^2.
 }
\]
Thus \(G\) is a strict counterexample exactly when
\[
 |\theta+1|<\sqrt{2k-2}
 \qquad(\theta\ne k).
\]
\end{theorem}

\begin{proof}
Girth at least five gives the distance-two matrix \(A_2=A^2-kI\), and diameter
three gives \(A_3=J-I-A-A_2\).  Substitute in
\(D=A+2A_2+3A_3\).  On \(\one^\perp\), \(J\) vanishes, and regularity gives
\(\delta^*=k\).  The principal distance eigenvalue is positive and is the
Perron root because every off-diagonal entry of \(D\) is positive.  The exact
operator and score identities are independently checked by
\codefile{scripts/verify_regular_score_calculus.py}.
\end{proof}

\begin{corollary}[Bipartite obstruction]\label{cor:bipartite}
A connected \(k\)-regular bipartite graph of girth at least five and diameter
three is not a strict counterexample for \(k\ge3\).
\end{corollary}

\begin{proof}
The nonprincipal eigenvalue \(-k\) satisfies
\[
 (-k+1)^2-(2k-2)=(k-1)(k-3)\ge0
 \qquad(k\ge3).
\]
Thus the required strict inequality fails, with equality exactly at \(k=3\).
\end{proof}

\begin{proposition}[Higher-diameter transfer]\label{prop:higher-transfer}
Let \(G\) be connected and \(k\)-regular, with diameter \(d\) and girth at
least \(2d-1\).  Let \(A_i\) denote the distance-\(i\) matrix, and define
\[
 F_0=1,
 \quad F_1=x,
 \quad F_2=x^2-k,
 \quad F_i=xF_{i-1}-(k-1)F_{i-2}\quad(i\ge3).
\]
Then \(A_i=F_i(A)\) for \(0\le i\le d-1\), and
\[
 D=dJ+q_d(A),
 \qquad
 q_d(x)=\sum_{i=0}^{d-1}(i-d)F_i(x).
\]
In particular,
\[
 q_3(x)=k-3-2x-x^2,
\]
\[
 q_4(x)=-x^3-2x^2+(2k-4)x+2k-4.
\]
\end{proposition}

\begin{proof}
Up to length \(d-1\), the girth condition makes nonbacktracking walks between
two vertices unique exactly when their length is the graph distance.  Hence the
distance-\(i\) matrices are the nonbacktracking polynomials in \(A\); summing
\(D=\sum_{i=0}^d iA_i\) and eliminating \(A_d\) with
\(J=\sum_{i=0}^d A_i\) proves the
formula.  This lies within the established distance-polynomial framework
\cite{HowladerPanigrahi2022,Fiol2016}.
\end{proof}

\section{Degree and diameter obstructions}\label{sec:obstructions}

\begin{lemma}\label{lem:diam-rayleigh}
For every connected graph,
\[
 \lambda_{\min}(D(G))\le-\diam(G).
\]
\end{lemma}

\begin{proof}
For a diametral pair \(u,v\), the Rayleigh quotient of \(e_u-e_v\) is
\(-d_G(u,v)\).
\end{proof}

\begin{theorem}\label{thm:regular-degree-six}
Every connected regular strict counterexample to WOW-284 has degree at least
six.
\end{theorem}

\begin{proof}
We use the LP ceiling proved independently in
Theorem~\ref{thm:lp-ceiling}; that theorem does not depend on the present
degree reduction.  Let the degree be \(k\).
Lemma~\ref{lem:diam-rayleigh} and strictness give \(\diam(G)<k\).
If \(k\le2\), then \(|V(G)|\ge3\) and connectedness give
\(\diam(G)\ge2\ge k\), a contradiction.  For \(k=3\), the radius-two lower
bound and \(\diam(G)\le2\) force equality in the Moore bound.  Hence \(G\) is
a degree-three Moore graph, and Theorem~\ref{thm:moore-threshold} gives
\(\Phi(G)=0\), not \(\Phi(G)>0\).

For \(k=4\), diameter two would require a degree-four Moore graph, whose
adjacency multiplicities are nonintegral.  In diameter three, the exact LP
bound of Theorem~\ref{thm:lp-ceiling} gives \(n<19\), whereas the
radius-two ball has \(17\) vertices and diameter three requires at least one
more.  Hence \(n=18\).  The distance-three matrix is then a perfect matching.
Its \(-1\)-eigenspace \(W\) is a nine-dimensional rational subspace of
\(\one^\perp\), and
\[
 A_3=J+3I-A-A^2
 \quad\Longrightarrow\quad
 (A^2+A-4I)|_W=0.
\]
The polynomial \(x^2+x-4\) is irreducible over \(\mathbb Q\), so a rational
space on which it annihilates an operator has even dimension, a contradiction.

For \(k=5\), diameter two again fails the Moore multiplicity condition.
A diametral geodesic in diameter four yields the principal submatrix
\(D(P_5)\), whose factor \(x^2+6x+4\) supplies the eigenvalue
\(-3-\sqrt5<-5\); Cauchy interlacing excludes this case.  In diameter
three, Meringer's lower bound and Theorem~\ref{thm:lp-ceiling} leave
\(n\in\{30,31,32\}\).  Since \(5n=2|E(G)|\), one has
\(n\equiv0\pmod2\), so \(n\in\{30,32\}\).  At \(n=32\), write
\(a\) for the average internal degree of the distance-two layer.
This layer has \(20\) vertices, each with one neighbour in the first layer,
so the number of edges from the second to the six-vertex third layer is
\(20(4-a)\le6\cdot5\); hence \(a\ge5/2\).  In the symmetric normalized layer
compression, the derivative of the only \(a\)-dependent block is
\[
 \begin{pmatrix}1&-\sqrt r\\-\sqrt r&r\end{pmatrix}
 =
 \begin{pmatrix}1\\-\sqrt r\end{pmatrix}
 \begin{pmatrix}1&-\sqrt r\end{pmatrix}\succeq0,
 \qquad r=\frac{5\cdot4}{6}.
\]
Thus every ordered compression eigenvalue is nondecreasing in \(a\).
At the smallest feasible value \(a=5/2\), the relevant nonprincipal
factor is
\[
 p_{5,6}(x)=4x^3+10x^2-16x-30,
 \qquad p_{5,6}(11/6)=-29/27<0.
\]
Because the leading coefficient is positive, its largest root exceeds
\(11/6>-1+\sqrt8\), contradicting the necessary bound
\(\theta<-1+\sqrt8\).  At \(n=30\), Meringer's isomorph-free enumeration
leaves exactly four \((5,5)\)-cages \cite{Meringer1999}; each fixed record has
an exact distance eigenvalue at most \(-5\).
The complete case split, the four fixed graph6 records, characteristic
polynomials, and interlacing directions are independently checked by
\codefile{scripts/verify_proof_audit_04_regular_low_degree.py}.
\end{proof}

\begin{theorem}[Endpoint-neighbourhood obstruction]\label{thm:endpoint-diameter}
Let \(G\) be any connected finite simple graph, and let \(u,v\) be vertices at
distance \(\ell=d_G(u,v)\ge5\).  Put \(p=d(u)\) and \(q=d(v)\).  Then
\[
 \resultbox{
 \lambda_{\min}(D(G))
 \le p+q-2-\sqrt{(p-q)^2+pq(\ell-2)^2}.
 }
\]
If \(\delta\) is the ordinary minimum degree, then
\[
 \resultbox{
 \lambda_{\min}(D(G))\le-\delta(\ell-4)-2.
 }
\]
Consequently every strict WOW-284 counterexample satisfies
\[
 \Delta>\delta(\ell-4)+2,
\]
where \(\Delta\) is the ordinary maximum degree.
In particular, every regular strict counterexample has diameter at most four.
\end{theorem}

\begin{proof}
The two endpoint neighbourhoods are disjoint.  Give weight \(a>0\) to
\(N(u)\), weight \(-b<0\) to \(N(v)\), and zero elsewhere.  Within one
neighbourhood distances are at most two; between the two neighbourhoods they
are at least \(\ell-2\).  Since the cross products are negative, the Rayleigh
quotient is at most that of
\[
 \begin{pmatrix}
 2(p-1)&-(\ell-2)\sqrt{pq}\\
 -(\ell-2)\sqrt{pq}&2(q-1)
 \end{pmatrix}.
\]
Its least eigenvalue is the first displayed bound, and its least eigenvector can
be chosen with both coordinates positive.  Write
\(p=\delta+\alpha\), \(q=\delta+\beta\), and \(t=\ell-2\).  The identity
\[
 (p-q)^2+pqt^2-(p+q+\delta(t-2))^2
 =(t-2)\{\delta t(\alpha+\beta)+(t+2)\alpha\beta\}
\]
is nonnegative and gives the second bound.  Finally \(\delta^*(G)\le\Delta\).
The sign choice, radical comparison, and integer rounding are independently
audited by \codefile{scripts/verify_proof_audit_10_endpoint_diameter.py}.
\end{proof}

\begin{theorem}[Diameter four]\label{thm:diameter-four}
Let \(G\) be connected, \(k\)-regular, of girth at least five and diameter
four.  Then
\[
 \resultbox{
 \lambda_{\min}(D(G))\le-\frac{7+\sqrt{16k+1}}2.
 }
\]
Consequently no such graph of degree \(2\le k\le9\) is a strict
counterexample.
\end{theorem}

\begin{proof}
Choose \(u,v\) at distance four and put \(U=N(u)\), \(V=N(v)\).  For fixed
\(a\in U\), distinct vertices \(b,b'\in V\) at distance two from \(a\)
cannot use the same common neighbour, or a \(4\)-cycle results.  Thus at most
\(k(k-1)\) pairs in \(U\times V\) have distance two, and
\[
 \sum_{a\in U,b\in V}d_G(a,b)\ge2k^2+k.
\]
Assign weights \(\alpha,\beta,-\alpha,-\beta\) to
\(u,U,v,V\), respectively.  Counting unordered pairs and then doubling gives
\[
 \frac{x^{\mathsf T}D(G)x}{x^{\mathsf T}x}
 \le
 \frac{-4\alpha^2-4k\alpha\beta-3k\beta^2}
 {\alpha^2+k\beta^2}.
\]
After setting \(y_1=\alpha\) and \(y_2=\sqrt{k}\,\beta\), the
right-hand side is the Rayleigh quotient of
\[
 \begin{pmatrix}-4&-2\sqrt{k}\\-2\sqrt{k}&-3\end{pmatrix},
\]
whose least eigenvalue is the displayed value and has a positive-coordinate
minimizer.  The strict comparison with \(-k\) holds for \(2\le k\le9\).
Every orientation factor, cross-distance sign, and endpoint comparison is
independently checked by
\codefile{scripts/verify_proof_audit_11_diameter_four.py}.
\end{proof}

\begin{corollary}[Regular trichotomy]\label{cor:trichotomy}
Every regular strict counterexample has one of the following forms:
\begin{enumerate}
\item diameter two, hence a Moore graph;
\item diameter three, with
\(\lvert\theta+1\rvert<\sqrt{2k-2}\) for every nonprincipal adjacency
eigenvalue \(\theta\);
\item diameter four, with degree at least ten.
\end{enumerate}
There are no regular strict counterexamples of diameter at least five.
\end{corollary}

\section{Moment bounds and the exact LP ceiling}\label{sec:lp}

Let \(G\) satisfy the hypotheses of Theorem~\ref{thm:diameter-three-score}, and
write its nonprincipal adjacency eigenvalues as \(\theta_i\).  Put
\(y_i=\theta_i+1\).

\begin{proposition}[Fourth-moment identity]\label{prop:moment-bound}
One has
\[
 \sum_{i=1}^{n-1}(2k-2-y_i^2)(y_i+1)^2
 =(k+2)\bigl((k+2)(k^2+3)-6n\bigr).
\]
Every strict counterexample therefore satisfies
\[
 \resultbox{
 n< B_k:=\frac{(k+2)(k^2+3)}6.
 }
\]
\end{proposition}

\begin{proof}
Use
\[
 \tr A=\tr A^3=0,
 \quad \tr A^2=nk,
 \quad \tr A^4=nk(2k-1),
\]
remove the principal eigenvalue, and expand both sides.  In a strict
counterexample, each factor \(2k-2-y_i^2\) is positive.  The sum cannot
vanish: otherwise every nonprincipal adjacency eigenvalue would equal
\(-2\), and \(\tr A=0\) would give
\(k-2(n-1)=0\), or \(n=(k+2)/2\), incompatible with the elementary
bound \(n\ge k+1\) for a simple \(k\)-regular graph.  The identity is
checked symbolically by
\codefile{scripts/verify_degree_six_gate.py} and independently within
\codefile{scripts/verify_proof_audit_02_two_sided_lp.py}.
\end{proof}

The preceding bound is in fact optimal within the entire standard one-variable
nonbacktracking LP hierarchy.

\begin{definition}
Let \(F_i=F_i^{(k)}\) be the nonbacktracking polynomials from
Proposition~\ref{prop:higher-transfer}, and set
\[
 I_k=[-1-\sqrt{2k-2},-1+\sqrt{2k-2}].
\]
A finite polynomial \(f=\sum_i f_iF_i\) is \emph{admissible} if
\[
 f_0>0,
 \qquad f_i\ge0\quad(i\ge5),
 \qquad f(x)\le0\quad(x\in I_k).
\]
\end{definition}

For a \(k\)-regular graph of girth at least five whose nonprincipal spectrum
lies in \(I_k\), one has \(\tr F_i(A)=0\) for \(1\le i\le4\), while
\(\tr F_i(A)\ge0\) for \(i\ge5\), since these traces count closed
nonbacktracking walks.  The coefficient conditions therefore give
\(nf_0\le\tr f(A)\).  On the other hand, \(f(\theta)\le0\) for every
nonprincipal eigenvalue, so \(\tr f(A)\le f(k)\).  Thus
\[
 nf_0\le\tr f(A)\le f(k),
 \qquad n\le\frac{f(k)}{f_0}.
\]

\begin{theorem}[Exact LP ceiling and rigidity]\label{thm:lp-ceiling}
For every integer \(k\ge4\) and every admissible \(f\),
\[
 \resultbox{
 \frac{f(k)}{f_0}\ge B_k=\frac{(k+2)(k^2+3)}6.
 }
\]
Equality holds if and only if \(f\) is a positive scalar multiple of
\[
 \resultbox{
 f_*(x)=\frac{(x+2)^2(x^2+2x-(2k-3))}{6(k+2)}.
 }
\]
Thus increasing the polynomial degree cannot improve this one-variable LP
bound.  Consequently, any connected \(k\)-regular graph of girth at least
five whose nonprincipal spectrum lies in the interior of \(I_k\) satisfies
\(n<B_k\).
\end{theorem}

\begin{proof}
The primal expansion is
\[
\begin{aligned}
6(k+2)f_*(x)={}&6(k+2)F_0(x)+2(2k+7)F_1(x)\\
&+(k+13)F_2(x)+6F_3(x)+F_4(x).
\end{aligned}
\]
On \(I_k\), the factor \((x+1)^2-(2k-2)\) is nonpositive, while
\((x+2)^2\ge0\); hence \(f_*\) is admissible and \(f_*(k)=B_k\).

For the dual certificate, put \(\Delta=\sqrt{2k-2}\),
\(\xi_\pm=-1\pm\Delta\), and \(\xi_0=-2\).  Define
\begin{align*}
w_-&=\frac{k(k+2)(2k^2-6-3(k-1)\Delta)}{24(2k-3)},\\
w_0&=\frac{k(k-1)(k^2+3)}{6(2k-3)},\\
w_+&=\frac{k(k+2)(2k^2-6+3(k-1)\Delta)}{24(2k-3)}.
\end{align*}
\(w_-,w_0,w_+>0\).  For the only nontrivial inequality, this follows from
\[
 (2k^2-6)^2-18(k-1)^3
 =2(k-3)(2k-3)(k^2+3)>0.
\]
The measure
\(\mu=w_-\delta_{\xi_-}+w_0\delta_{\xi_0}+w_+\delta_{\xi_+}\) satisfies
\[
 \mu(1)=B_k-1,
 \qquad \mu(F_i)=-F_i(k)\quad(1\le i\le4).
\]
For \(i\ge5\), the slack \(a_i=\mu(F_i)+F_i(k)\) is strictly positive.  For
\(5\le i\le9\), exact calculation gives, after removing the common positive
factor \(k(k-1)(k+2)(k^2+3)/6\), respectively,
\[
 2,\quad 5k-13,\quad 2(3k^2-17k+25),
\]
\[
 6k^3-47k^2+139k-150,
 \quad
 2(3k^4-27k^3+106k^2-219k+194),
\]
all positive for \(k\ge4\); for the nontrivial residual factors this follows
after writing \(k=m+4\), when all coefficients are nonnegative and the
constant terms are positive.  For \(i\ge10\), put \(r=k-1\ge3\).
The support lies in \([-2\sqrt r,2\sqrt r]\), because
\(1+\sqrt{2r}\le2\sqrt r\).  For \(|z|\le1\), the recurrence gives
\[
 F_i(2\sqrt r\,z)
 =r^{i/2}U_i(z)-r^{(i-2)/2}U_{i-2}(z),
\]
where \(U_j\) is the Chebyshev polynomial of the second kind.  Using
\(|U_j(z)|\le j+1\) yields
\[
 \frac{|\mu(F_i)|}{F_i(k)}
 \le \frac{2i+1}{3}\,3^{3-i/2}.
\]
At \(i=10\) the right-hand side is \(7/9\), and it decreases thereafter
because \(3(2i+1)^2-(2i+3)^2=8i^2-6>0\).  Hence
\[
 \frac{|\mu(F_i)|}{F_i(k)}<1\qquad(i\ge10).
\]
Expanding in the nonbacktracking basis gives
\[
 \int f\,d\mu
 =B_kf_0-f(k)+\sum_{i\ge5}f_i a_i
 \ge B_kf_0-f(k).
\]
Since \(f\le0\) on \(\operatorname{supp}\mu\),
\(\int f\,d\mu\le0\), which proves weak duality.
If equality holds, strict positivity of every high-degree slack forces
\(f_i=0\) for \(i\ge5\).  Equality on the positive dual support forces zeros
at \(\xi_-,-2,\xi_+\); the interior zero \(-2\) has even multiplicity because
\(f\le0\) on \(I_k\).  Degree at most four then forces \(f\) to be a scalar
multiple of \(f_*\), and \(f_0>0\) makes the scalar positive.  Equality in
the graph bound would then force every nonprincipal adjacency eigenvalue to be
\(-2\), which contradicts the trace equation; hence the open-window bound is
strict.  Every symbolic identity, finite slack, tail bound, and
equality-nullspace calculation is independently checked by
\codefile{scripts/verify_proof_audit_02_two_sided_lp.py}.
\end{proof}

\begin{theorem}[Integral optimal-slack bound]\label{thm:integral-slack}
Let \(G\) be connected and \(k\)-regular, of girth at least five and diameter
three, and suppose
\[
 \lvert\theta+1\rvert<\sqrt{2k-2}
 \qquad(\theta\ne k)
\]
for every adjacency eigenvalue \(\theta\).  Define
\[
 g_k(x)=(x+2)^2\bigl((x+1)^2-(2k-2)\bigr),\quad
 C_k=(k+2)^2(k^2+3),\quad h_k=6(k+2),
\]
\[
 \mathcal S_k=-g_k(A)+\frac{C_k}{n}J.
\]
Then
\[
 \mathcal S_k\succeq0,\qquad
 \mathcal S_k\one=0,\qquad
 \tr\mathcal S_k=h_k(B_k-n).
\]
Moreover,
\[
 \mathcal E_k=g_k(A)-(h_k+1)J+I
\]
is a nonzero, symmetric, entrywise nonnegative integral matrix with zero
diagonal and constant row sum
\[
 \varepsilon_{k,n}=C_k-(h_k+1)n+1.
\]
On \(\one^\perp\), one has \(\mathcal S_k=I-\mathcal E_k\).  In particular,
\[
 \resultbox{
 n\le
 \left\lfloor
 \frac{(k+2)^2(k^2+3)}{6(k+2)+1}
 \right\rfloor.
 }
\]
\end{theorem}

\begin{proof}
The polynomial \(g_k=6(k+2)f_*\) is the optimal LP polynomial from
Theorem~\ref{thm:lp-ceiling}.  Thus \(\mathcal S_k\) vanishes on
\(\langle\one\rangle\), while on a nonprincipal \(\theta\)-eigenspace its
eigenvalue is
\[
 \bigl(2k-2-(\theta+1)^2\bigr)(\theta+2)^2\ge0.
\]
The nonbacktracking expansion of \(g_k\) has constant coefficient \(h_k\).
Girth at least five gives
\((F_i(A))_{uu}=0\) for every \(u\) and \(1\le i\le4\), and hence
\[
 \tr\mathcal S_k=C_k-h_kn=h_k(B_k-n),
 \qquad (g_k(A))_{uu}=h_k.
\]

Put \(z_{uv}=(g_k(A))_{uv}\in\mathbb Z\) for \(u\ne v\), and set
\(a=C_k/n-h_k=(\mathcal S_k)_{uu}\).  The \(2\times2\) principal minor on
\(\{u,v\}\) gives
\[
 \left|\frac{C_k}{n}-z_{uv}\right|\le a.
\]
Equality on the upper side would give
\((e_u-e_v)^{\mathsf T}\mathcal S_k(e_u-e_v)=0\), hence
\(e_u-e_v\in\ker\mathcal S_k\).  The strict spectral window gives
\[
 \ker\mathcal S_k=\langle\one\rangle\oplus E_{-2}(A).
\]
Since \(e_u-e_v\perp\one\), it would be a \((-2)\)-eigenvector of \(A\).
Yet the \(u\)-coordinate of \(A(e_u-e_v)\) is \(-1\) if \(u\sim v\) and
\(0\) otherwise, never \(-2\).  Therefore
\[
 z_{uv}\ge h_k+1.
\]
It follows directly that \(\mathcal E_k\) has the stated entrywise properties,
and \(g_k(A)\one=C_k\one\) gives its row sum.  The identity on
\(\one^\perp\) follows by eliminating \(J\).

It remains to prove \(\mathcal E_k\ne0\).  Otherwise
\[
 g_k(A)=(h_k+1)J-I,
\]
so \(p_k(A)=0\) on \(\one^\perp\), where \(p_k(x)=g_k(x)+1\).  After
\(y=x+2\),
\[
 p_k(y-2)=y^4-2y^3+(3-2k)y^2+1.
\]
This polynomial is irreducible over \(\mathbb Q\).  Its only possible rational
roots are \(\pm1\), whose values are \(3-2k\) and \(7-2k\).  A factorization
into monic integer quadratics would have constant terms both \(1\) or both
\(-1\); these alternatives force the cubic and linear coefficients to be,
respectively, equal or opposite, whereas they are \(-2\) and \(0\).

Rational canonical form now gives
\[
 \chi_{A|_{\one^\perp}}(x)=p_k(x)^m,\qquad n-1=4m.
\]
The four roots of \(p_k\) sum to \(-6\), so \(\tr A=0\) gives
\[
 0=k-6m.
\]
Hence \(n-1=4m=2k/3\), contradicting \(n\ge k+1\).  Thus
\(\mathcal E_k\ne0\).  Its constant row sum is consequently a positive
integer, so \(\varepsilon_{k,n}\ge1\), which is equivalent to the displayed
order bound.  The symbolic expansion, irreducibility alternatives, and finite
specializations have also been checked by independent exact audits.
\end{proof}

\section{Optimal-slack integrality and local positivity}
\label{sec:degree-six}

The trace of \(\mathcal S_k\) is the one-variable LP defect.  Its
\(2\times2\) principal minors give integral local restrictions, and larger
principal minors form a canonical semidefinite hierarchy.  We retain the
edge-local argument below because it exposes the cycle-count obstruction
hidden by the stronger global order bound.

\begin{theorem}[Three-to-one excess bound]\label{thm:three-to-one}
Under the hypotheses of Theorem~\ref{thm:integral-slack}, assume \(k\ge6\)
and write \(\varepsilon=\varepsilon_{k,n}\).  The integral parameter
\[
 r=2\varepsilon-n-2
   =2(k+2)^2(k^2+3)-(12k+27)n
\]
satisfies
\[
 r>0,\qquad n\le3r.
\]
Consequently
\[
 \resultbox{
 n\le\left\lfloor
 \frac{3(k+2)^2(k^2+3)}{18k+41}
 \right\rfloor.
 }
\]
\end{theorem}

\begin{proof}
Put \(E=\mathcal E_k\), \(C=C_k\), \(h=h_k\), and
\[
 \rho=\frac{\varepsilon-1}{n}
      =\frac{1+r/n}{2}.
\]
The optimal slack matrix has the form
\[
 \mathcal S_k=I-E+\rho J\succeq0.
\]
 We first record the divisibility identity
\[
 128(2C-r)
 =(4k+9)(64k^3+112k^2+196k+327)+(129-128r).
\]
Since \(n=(2C-r)/(12k+27)\) is integral and \(4k+9\) is odd,
\begin{equation}\label{eq:r-divisibility}
 4k+9\mid129-128r.
\end{equation}

\smallskip
\noindent\emph{Positivity of \(r\).}
Suppose \(r\le0\).  A \(2\times2\) principal minor of \(\mathcal S_k\)
gives
\[
 E_{uv}\le2+\frac rn\qquad(u\ne v).
\]
Thus \(E\) is the adjacency matrix of a simple graph when \(r<0\).  If
\(r=0\) and \(E_{uv}=2\), positivity puts \(e_u+e_v\) in
\(\ker\mathcal S_k\).  After subtracting its projection onto \(\one\), the
strict spectral window gives a \((-2)\)-eigenvector of \(A\).  Its
\(u\)-coordinate would require
\[
 A_{uv}=-2+\frac{2k+4}{n},
\]
which is impossible because \(n\ge k^2+2\) and \(A_{uv}\in\{0,1\}\).
Hence \(E\) is simple also when \(r=0\).

Let \(X=\overline E\).  It is regular of degree
\[
 d=n-1-\varepsilon=\frac{n-r-4}{2},
\]
and \(A(X)=-I-E+J\) has least eigenvalue at least \(-2\).  Every component
has at least \(d+1\) vertices, so \(X\) has at most two components.  We use
the fact that \(E\), being a polynomial in \(A\) and \(J\), commutes with
\(A\); hence every rational eigenspace of \(E\) is \(A\)-invariant.  We use
 the classification of connected regular graphs of order greater than \(28\)
 and least eigenvalue at least \(-2\): such a graph is a line graph or a
 cocktail-party graph \cite[Theorem~4.1.1]{CvetkovicEtAl2004};
 see also \cite{KoolenEtAl2025}.

Suppose first that \(X\) is connected.  The complete case would give
\(E=0\), excluded in Theorem~\ref{thm:integral-slack}.  In the
cocktail-party case, \(E\) is a perfect matching and \(r=-n\).  The identity
\[
 1296C=(6k+13)(216k^3+396k^2+654k+1175)+277
\]
forces \(k=44\) and \(n=14812\).  On the \(7406\)-dimensional
 \((-1)\)-eigenspace of \(E\), the rational operator \(A\) is annihilated by
 \(g_{44}(x)+2\).

If \(X=L(Y)\) and \(Y\) is \(q\)-regular, then \(q\ge n/4\) and
\(|V(Y)|\le8\), contradicting \(q\le |V(Y)|-1\) because \(n\ge38\).
If \(Y\) is semiregular bipartite with part sizes \(a\ge b\ge2\), then
\[
 \frac1a+\frac1b=\frac{n-r}{2n}\ge\frac12,\qquad n\le ab.
\]
The cases \(b\ge3\) have \(ab<38\).  For \(b=2\), connectedness forces
\(Y=K_{a,2}\), hence \(r=-4\).  Equation~\eqref{eq:r-divisibility} leaves
\(k=158\) and \(n=664748\).  The resulting \((-1)\)-eigenspace of \(E\)
 has dimension \(332373\).

If \(X\) has two components, write their orders as
\(d+1+a_1\) and \(d+1+a_2\).  Then
\[
 a_1+a_2=r+2,
\]
so \(r\in\{-2,-1,0\}\).  Checking the divisors in
\eqref{eq:r-divisibility}, together with \(n\ge k^2+2\) and \(kn\) even,
leaves only
\[
 (r,k,n)=(-1,62,40875).
\]
The components then have orders \(20437\) and \(20438\); one is complete
and the other cocktail-party.  Thus \(E\) has eigenvalue \(-1\) with
 multiplicity \(10219\).

 In all three exceptional cases the rational primary space is annihilated by
 \(g_k(x)+2\).  This quartic is irreducible over \(\mathbb Q\) for every
 integer \(k\ge6\), except \(k=7\).  Indeed, after \(y=x+2\) it becomes
 \[
  y^4-2y^3+(3-2k)y^2+2.
 \]
 Gauss's lemma reduces a quadratic factorization to monic integer quadratics
 whose constant terms multiply to \(2\); coefficient comparison leaves only
 \(k=4\) or \(k=7\), while the rational-root alternatives give \(k=2\) or
 \(k=4\).  The degrees \(44,62,158\) are therefore irreducible cases, but the
 corresponding dimensions \(7406,10219,332373\) are not divisible by four.
 This contradiction proves \(r>0\).

 \smallskip
\noindent\emph{Simplicity above the putative boundary.}
 Assume for contradiction that \(n>3r\), and put \(x=r/n\in(0,1/3)\).
 The \(2\times2\) minors give \(E_{uv}\le2\).  Suppose \(E_{uv}=2\), and
 for \(w\notin\{u,v\}\) set \(s_w=E_{uw}+E_{vw}\).  Cauchy--Schwarz for
 \(e_u+e_v\) and \(e_w\) in the Gram matrix \(\mathcal S_k\) gives
\[
 (1+x-s_w)^2\le x(3+x),
\]
 so the nonnegative integer \(s_w\) belongs to \(\{1,2\}\).  The row sums
 show that exactly \(r\) vertices have \(s_w=2\).  Let \(W\) be their set,
 let \(y=\sum_{w\in W}e_w\), and put
\[
  e_W=\sum_{\{w,z\}\subseteq W}E_{wz}.
\]
 Then
\[
 \begin{aligned}
  (e_u+e_v)^{\mathsf T}\mathcal S_k(e_u+e_v)&=2x,\\
  (e_u+e_v)^{\mathsf T}\mathcal S_k y&=r(x-1),\\
  y^{\mathsf T}\mathcal S_k y&=r+\frac{r^2(1+x)}2-2e_W.
 \end{aligned}
\]
 Positivity of this \(2\times2\) Gram determinant yields
 \[
  2xr+r^2(3x-1)-4xe_W\ge0.
 \]
 Since \(e_W\ge0\), it follows that \(n\le3r+2\).  Hence
 \(n=3r+t\) with \(t\in\{1,2\}\).  Substitution in the definition of \(r\)
 shows that \(18k+41\) divides a fixed remainder:
 \[
 \begin{array}{c|c|c}
  t&\text{remainder}&\text{admissible }(18k+41,k)\\
  \hline
  1&167642&\text{none}\\
  2&202634&(1427,77).
 \end{array}
 \]
 The remaining case gives \((k,n,r)=(77,77831,25943)\), but \(kn\) is odd,
 contrary to the handshake lemma.  Therefore \(E\) is simple.

 Again \(X=\overline E\) is \(d\)-regular with least eigenvalue at least
 \(-2\).  If \(X\) had at least three components, then
 \(n\ge3(d+1)\), or \(n\le3r+6\).  Thus \(n=3r+t\) for
 \(1\le t\le6\).  The same fixed-remainder calculation gives
 \[
 \begin{array}{c|rrrrrr}
  t&1&2&3&4&5&6\\ \hline
  \text{remainder}
   &167642&202634&237626&272618&307610&342602,
 \end{array}
 \]
 and leaves no integral graph: the \(t=2\) case is excluded above, while the
 sole divisor candidate for \(t=3\) does not make \(r\) integral.  Hence \(X\)
 has at most two components.

 Suppose first that \(X\) is connected.  The cocktail-party case contradicts
 \(r>0\), so \(X=L(Y)\).  If \(Y\) is \(q\)-regular on \(v\) vertices, then
 \[
  q=\frac{n-r}{4}>\frac n6,\qquad v=\frac{2n}{q}<12.
 \]
 Simplicity and \(n\ge38\) leave only
 \[
 \begin{aligned}
  (q,v;n,r)\in\{&(8,10;40,8),(9,10;45,9),\\
                 &(8,11;44,12),(10,11;55,15)\}.
 \end{aligned}
 \]
 The radius-two lower bound forces \(k\le7\), and direct substitution in the
 defining formula for \(r\) excludes all four cases.

 For a semiregular bipartite root with part sizes \(a\ge b\ge2\),
 write \(p=n/a\le q=n/b\) for its two degrees.  Then
 \[
  p+q=\frac{n-r}{2}>\frac n3,
  \qquad (b-2)n=b(r+2p),
 \]
 so \(b<6\).  Connectedness excludes \(p=1\).  The cases \(b=2\) and
 \(b=5\) are immediate; \(b=4\) leaves only
 \[
  (p,q,a,b;n,r)=(4,10,10,4;40,12),\ (4,11,11,4;44,14),
 \]
 both excluded by the radius-two bound and the formula for \(r\).
 For \(b=3\), one has \(p=2\) or \(3\), giving \(n=3r+12\) or
 \(n=3r+18\).  The corresponding fixed remainders are
 \(552554=2\cdot276277\) and \(762506=2\cdot381253\).  Their odd cofactors
 are prime and congruent to \(13\pmod {18}\), whereas \(18k+41\equiv5
 \pmod {18}\).  Thus \(X\) cannot be connected.

 It remains to consider two components.  For \(k=6,7,8\), the inequalities
 \(r>0\) and \(n>3r\) contain no admissible integral order.  For \(k\ge9\),
 one has \(n>150\), and both components have order greater than \(28\).
 A regular line-graph root is too small.  For a semiregular bipartite root
 of a component of order \(N\), the part sizes satisfy
 \[
  \frac1a+\frac1b=\frac{d+2}{N}
  \ge\frac{n-r}{n+r+2}>\frac{49}{100}.
 \]
 If \(b\ge3\), the five possible pairs \((a,b)\) have \(ab\le18<N\).
 If \(b=2\), connectedness forces the root \(K_{d,2}\), whose line graph has
 order \(2d\).  Thus every component is \(K_{d+1}\), a cocktail-party graph
 of order \(d+2\), or \(L(K_{d,2})\).  The first two types alone force
 \(r\le0\).  If exactly one component has order \(2d\), then
 \(n=3r+8\) or \(n=3r+10\); if both do, then \(n=2r+8\).
 The respective fixed remainders are
 \[
  412586,\qquad482570,\qquad1792898.
 \]
 The first and third have no admissible divisor of the required linear form.
 The second leaves only \(k=123\), for which \(kn\) is odd.  This final
 contradiction proves \(n\le3r\).

 Substituting the definition of \(r\) gives
 \[
  (18k+41)n\le3(k+2)^2(k^2+3),
 \]
 which is the displayed bound.  All polynomial divisions, irreducibility
 alternatives, fixed-remainder cases, line-root reductions, and Gram
 determinants in this argument have been replayed independently in exact
 arithmetic.
\end{proof}

\begin{corollary}[Rigidity at equality]\label{cor:three-to-one-equality}
Under the hypotheses and notation of Theorem~\ref{thm:three-to-one}, equality
in \(n\le3r\) can occur only for
\[
 \resultbox{(k,n,r)=(103,185220,61740).}
\]
This is an arithmetic parameter classification; it does not assert the
existence of a graph with these parameters.
\end{corollary}

\begin{proof}
Equality gives
\[
 (18k+41)r=(k+2)^2(k^2+3)=:C_k.
\]
Exact Euclidean division yields
\[
\begin{aligned}
18^4C_k={}&(18k+41)
 (5832k^3+10044k^2+17946k+29107)\\
&+66325,
\end{aligned}
\]
where \(66325=5^2\cdot7\cdot379\).  Hence \(18k+41\mid66325\).
Since \(k\ge6\), this divisor is at least \(149\) and is congruent to \(5\)
modulo \(18\).  Among the positive divisors of \(66325\), only
\(1895\) has these properties.  Thus \(k=103\), after which direct
substitution gives \(r=61740\) and \(n=3r=185220\).
\end{proof}

\begin{proposition}[Signed-complement bridge]\label{prop:signed-complement}
Under the hypotheses and notation of Theorem~\ref{thm:three-to-one}, assume
\(0<r<n\) and define
\[
 S=(6k+14)J-2I-g_k(A).
\]
Then \(S\) is a signed adjacency matrix:
\[
 S_{uu}=0,\qquad S_{uv}\in\{-1,0,1\}\quad(u\ne v).
\]
It has constant signed row sum \((n-r-4)/2\), and
\[
 S+2I\succeq0,\qquad E_{-2}(S)=E_{-2}(A).
\]
\end{proposition}

\begin{proof}
Write \(x=r/n\in(0,1)\), \(\rho=(1+x)/2\), and
\(E=\mathcal E_k\).  Since
\[
 \mathcal S_k=I-E+\rho J\succeq0,
\]
its \(2\times2\) principal minors give
\(-1\le E_{uv}\le2+x\).  The entries of \(E\) are nonnegative integers, so
\(E_{uv}\in\{0,1,2\}\) for \(u\ne v\), and \(S=J-I-E\) has the asserted
entries.  The row sum follows from
\(\varepsilon=(n+r+2)/2\).  Moreover,
\[
 S+2I
 =\mathcal S_k+\frac{n-r}{2n}J\succeq0.
\]
The matrices \(A\) and \(S\) commute.  On a nonprincipal adjacency
eigenvector with eigenvalue \(\theta\), the corresponding eigenvalue of \(S\)
is \(-2-g_k(\theta)\).  The only zero of \(g_k\) in the open shifted WOW
window is \(\theta=-2\); the other two zeros are its excluded endpoints.
The principal eigenvalue of \(S\) is not \(-2\) because \(r<n\).  Hence the
two \((-2)\)-eigenspaces coincide.
\end{proof}

\begin{proposition}[Edge-local cycle bounds]\label{prop:edge-cycle}
Let \(G\) be connected and \(k\)-regular, of girth at least five and diameter
three, and suppose
\[
 |\theta+1|\le\sqrt{2k-2}
 \qquad(\theta\ne k)
\]
for every adjacency eigenvalue \(\theta\).  For an edge \(uv\), let
\(\sigma_{uv}\) be the number of
\(5\)-cycles containing that edge.  Then
\[
 2k-2\le\sigma_{uv}\le
 \frac{2(k+2)^2(k^2+3)}n-10k-26.
\]
If \(n=k^2+1+c\), then also
\[
 \sigma_{uv}\ge(k-1)^2-c.
\]
\end{proposition}

\begin{proof}
Put
\[
 f_k(x)=(x+2)^2((x+1)^2-(2k-2)),
 \quad C_k=f_k(k),
\]
and
\[
 M=-f_k(A)+\frac{C_k}{n}J.
\]
The spectral window gives \(M\succeq0\).  For an edge \(uv\),
\[
 (A^3)_{uv}=\sum_{z\sim v}(A^2)_{uz}=k+(k-1)=2k-1.
\]
Here \(z=u\) contributes \(k\), while every other neighbour of \(v\) is at
distance two from \(u\) and has a unique length-two path from \(u\).
Moreover,
\[
 (A^4)_{uv}=\sum_z(A^2)_{uz}(A^2)_{zv}=\sigma_{uv}.
\]
Indeed, the nonzero summands away from the diagonal are precisely the vertices
at distance two from both \(u\) and \(v\).  Their two unique length-two
paths, together with \(uv\), form a five-cycle, and each five-cycle through
\(uv\) yields one such vertex.  Since
\(C_k=(k+2)^2(k^2+3)\), the diagonal and edge entries of \(M\) are
\[
 a=\frac{C_k}{n}-6(k+2),
 \qquad
 b=\frac{C_k}{n}-(4k+14)-\sigma_{uv}.
\]
The principal submatrix on \(\{u,v\}\) is
\(\bigl(\begin{smallmatrix}a&b\\b&a\end{smallmatrix}\bigr)\), so
\(a\ge0\) and \(-a\le b\le a\).  The inequality \(b\le a\) gives
\(\sigma_{uv}\ge2k-2\), while \(b\ge-a\) gives the stated upper bound.

For the final bound, every radius-two ball has size \(k^2+1\).  The set
\[
 \{u,v\}\cup(N(u)\setminus\{v\})\cup(N(v)\setminus\{u\})
\]
contains \(2k\) vertices and lies in \(B_2(u)\cap B_2(v)\).  Every further
intersection vertex is at distance two from both endpoints and is therefore in
the preceding five-cycle bijection.  Hence
\[
 |B_2(u)\cap B_2(v)|=2k+\sigma_{uv}.
\]
Inclusion--exclusion and \(n=k^2+1+c\) now give
\(\sigma_{uv}\ge(k-1)^2-c\).  The complete walk classification, sign
directions, and radius-two bijection have also been checked independently and
exactly.
\end{proof}

\begin{corollary}[Edge--cycle divisibility sieve]
\label{cor:edge-cycle-sieve}
Under the hypotheses of Proposition~\ref{prop:edge-cycle}, put
\[
 L_{k,n}=\max\{2k-2,\;2k^2-2k+2-n\},
 \qquad
 U_{k,n}=\frac{2(k+2)^2(k^2+3)}{n}-10k-26.
\]
Necessarily
\[
 \resultbox{\lceil L_{k,n}\rceil\le\lfloor U_{k,n}\rfloor.}
\]
If both sides equal an integer \(s\), then every edge of \(G\) lies in
exactly \(s\) five-cycles and
\[
 \resultbox{5\mid \frac{skn}{2}.}
\]
\end{corollary}

\begin{proof}
Writing \(n=k^2+1+c\), the two lower bounds in
Proposition~\ref{prop:edge-cycle} combine to
\[
 \sigma_{uv}\ge
 \max\{2k-2,(k-1)^2-c\}=L_{k,n},
\]
while its upper bound is \(U_{k,n}\).  Since \(\sigma_{uv}\) is an integer,
the first conclusion follows.  If the two integer bounds coincide at \(s\),
then \(\sigma_{uv}=s\) for every edge.  Counting edge--five-cycle incidences
gives
\[
 5N_5=\sum_{uv\in E(G)}\sigma_{uv}
 =s|E(G)|=\frac{skn}{2},
\]
which proves the divisibility condition.
\end{proof}

\begin{theorem}\label{thm:degree-six-fifty}
Every connected \(6\)-regular strict counterexample to WOW-284 has order at
most \(50\).
\end{theorem}

\begin{proof}
A separate Rayleigh and trace argument shows that any degree-six strict
counterexample has diameter three.  For vertices at distance \(d\ge4\), the
vector with weights \(3,1,-3,-1\) on the two endpoints and their respective
neighbourhoods has Rayleigh quotient at most
\[
 \frac{204-81d}{15}\le-8,
\]
contradicting \(\delta^*=6\).  Diameter two would force equality in the
Moore bound, hence order \(37\) and adjacency characteristic polynomial
\((x-6)(x^2+x-5)^{18}\); its root sum is \(-12\), contradicting
\(\tr A=0\).  Thus the diameter is three, and
 Theorem~\ref{thm:three-to-one} gives
\[
 n\le\left\lfloor\frac{3\cdot8^2\cdot39}{149}\right\rfloor=50.
\]

For an independent local explanation of the excluded boundary, assume
\(n=51\).  Corollary~\ref{cor:edge-cycle-sieve} has
\(\lceil L_{6,51}\rceil=\lfloor U_{6,51}\rfloor=11\), and hence requires
\[
 5N_5=153\cdot11=1683,
 \qquad 5\nmid 1683,
\]
which is impossible.  An independent exact audit verifies the full diameter
reduction and incidence calculation.
\end{proof}

\begin{corollary}[Low-degree order windows]\label{cor:low-degree-windows}
There is no regular strict counterexample of degree at most five.  In degrees
six through nine, every regular strict counterexample satisfies
\begin{align*}
k=6&:\quad n\le50,\\
k=7&:\quad n=50\text{ in diameter two, or }n\le74\text{ in diameter three},\\
k=8&:\quad n\le108,\\
k=9&:\quad n\le150.
\end{align*}
\end{corollary}

\begin{proof}
The degree exclusion is Theorem~\ref{thm:regular-degree-six}, and
Theorem~\ref{thm:diameter-four} removes diameter four for \(k\le9\).
 For diameter three, Theorem~\ref{thm:three-to-one} gives, for
\(k=6,7,8,9\), respectively,
\[
 n\le50,\ 75,\ 108,\ 150.
\]
Since \(7n=2|E(G)|\), the degree-seven bound improves to \(74\).  The
diameter-two alternatives are determined by the Moore multiplicities: only
the degree-seven, order-\(50\) Hoffman--Singleton case survives.
\end{proof}

For comparison, the unadjusted diameter-three bounds in degrees \(6\) through
\(20\) are
\[
\begin{array}{c|rrrrrrrrrrrrrrr}
 k&6&7&8&9&10&11&12&13&14&15&16&17&18&19&20\\
 \hline
 n&50&75&108&150&201&263&336&422&521&635&765&911&1075&1257&1459.
\end{array}
\]
When \(k\) is odd, \(kn=2|E(G)|\) requires \(n\) to be even; in particular,
the entries for \(k=7,11,15,17,19\) improve to
\(74,262,634,910,1256\), respectively.

\begin{proposition}[The order-\(50\) \((-2)\)-multiplicity bound]
\label{prop:order50-minus-two}
Let \(G\) be \(6\)-regular, of order \(50\), and of girth at least five.
The multiplicity \(m_{-2}(A)\) of the adjacency eigenvalue \(-2\) satisfies
\[
 m_{-2}(A)\le20.
\]
\end{proposition}

\begin{proof}
Put \(m=m_{-2}(A)\), remove the principal eigenvalue \(6\) and the \(m\)
copies of \(-2\), and denote the remaining spectral moments by
\(\mu_0,\ldots,\mu_4\).  Girth at least five gives
\[
 \tr A=0,\quad \tr A^2=300,\quad \tr A^3=0,\quad \tr A^4=3300,
\]
and therefore
\[
 (\mu_0,\mu_1,\mu_2,\mu_3,\mu_4)
 =(49-m,-6+2m,264-4m,-216+8m,2004-16m).
\]
The moment matrix
\[
 H=\begin{pmatrix}
 \mu_0&\mu_1&\mu_2\\
 \mu_1&\mu_2&\mu_3\\
 \mu_2&\mu_3&\mu_4
 \end{pmatrix}
\]
is positive semidefinite.  Exact expansion gives
\[
 \det H=3600(1625-81m)\ge0.
\]
Since \(m\) is an integer, \(m\le20\).
\end{proof}

\subsection{Necessary structure at order fifty}

The remaining degree-six boundary is highly constrained, although not yet
eliminated.

\begin{theorem}\label{thm:order50-feasibility}
Let \(G\) be a connected \(6\)-regular graph of order \(50\) and girth
at least five, and suppose \(\Phi(G)>0\).  Then \(G\) has diameter three.
Every edge lies in \(12\) or \(13\) five-cycles.  Let \(H\) be the spanning
subgraph of edges of the second type, let \(m=|E(H)|\), and let \(\tau(v)\) be
the number of five-cycles through \(v\).  Then
\[
 \tau(v)\in\{36,37,38\},
 \qquad d_H(v)=2\tau(v)-72\in\{0,2,4\},
\]
\[
 m\equiv0\pmod5,
 \qquad N_5=360+\frac m5.
\]
For a two-edge path \(u-v-w\), put
\[
 r_{uvw}=6\alpha_{uvw}+\beta_{uvw},
\]
where \(\alpha\) and \(\beta\) count the five- and six-cycles containing the
path.  The allowed values are
\[
\begin{array}{c|c}
\text{types of the two incident edges}&r_{uvw}\\
\hline
\text{low--low}&30,31,32\\
\text{mixed}&30,31,32\\
\text{high--high}&30,31.
\end{array}
\]
Writing \(S_2=\sum_vd_H(v)^2\), one has
\[
 1950-m\le N_6
 \le2200-\frac{5m}{6}-\frac{S_2}{12},
\]
\[
 \frac{43m^2-70200m+119632500}{58500}
 \le N_6
 \le\frac{4220000-2200m-7m^2}{2000}.
\]
Exact enumeration of the resulting coarse degree profiles leaves \(266\)
possibilities.
\end{theorem}

\begin{proof}
The diameter reduction used in Theorem~\ref{thm:degree-six-fifty} applies
verbatim: diameter at least four gives a Rayleigh quotient at most \(-8\), and
diameter two gives the trace contradiction from
\((x-6)(x^2+x-5)^{18}\).  Hence the diameter is three.  Around a fixed vertex
the distance layers have sizes \(1,6,30,13\).  If \(\tau\) is the number of
five-cycles through the centre, the average row quotient is similar to the
symmetric compression on normalized layer indicators.  Its nonprincipal factor
\(q_\tau\) satisfies
\[
 195q_\tau(-1+\sqrt{10})
 =(-215+56\sqrt{10})\tau+7350-1860\sqrt{10}.
\]
The coefficient of \(\tau\) is negative because
\(56^2\cdot10<215^2\), and at \(\tau=39\) the right-hand side is
\(9(-115+36\sqrt{10})<0\).  Also
\[
 195q_\tau(6)=1500(75-\tau)>0.
\]
The last inequality is strict: the third layer is nonempty and connected to the
second, while their edge count is \(150-2\tau\).  Thus \(\tau\ge39\)
would place a nonprincipal compression eigenvalue in
\((-1+\sqrt{10},6)\), contradicting interlacing and the open WOW window.
Conversely, \(150-2\tau\le13\cdot6\), so \(\tau\ge36\).  Hence
\(\tau\in\{36,37,38\}\).  Finally,
\(\sum_{e\ni v}\sigma_e=2\tau(v)\), so the high-edge degree is
\(2\tau(v)-72\); counting edge--five-cycle incidences gives
\(5N_5=12\cdot150+m\).

For a two-path \(u-v-w\), girth at least five gives
\((A^3)_{uw}=\alpha_{uvw}\) and
\((A^4)_{uw}=16+\beta_{uvw}\).  The corresponding \(3\times3\) principal
minor of the centered positive-semidefinite matrix yields the displayed finite
sets for \(r_{uvw}\), except for an apparent equality value \(r=29\).
At equality, the Gram norm of \(e_u-e_w\) vanishes, so this vector lies in
the kernel of the centered matrix.  On \(\one^\perp\), that matrix is
\(-g_6(A)\).  The strict shifted window excludes the two endpoint zeros of
\(g_6\), leaving only its double zero at \(-2\); hence the kernel there is
exactly the adjacency \(-2\)-eigenspace.  However, the \(u\)-coordinate of
\(A(e_u-e_w)\) is zero because \(u\not\sim w\), whereas the
\(u\)-coordinate of \(-2(e_u-e_w)\) is \(-2\).  This excludes
\(r=29\).  Summing the local inequalities gives the first pair of
\(N_6\) bounds; shifted moment and localizing matrices give the second.
The complete symbolic determinants, Schur complements, kernel argument, and
integer enumeration have also been checked independently and exactly.  The
surviving \(266\) profiles are exact pruning data, not an existence or
nonexistence claim.
\end{proof}

\begin{remark}[Signed-root Gram formulation]\label{rem:signed-root}
Under the hypotheses of Theorem~\ref{thm:order50-feasibility}, put
\[
 T=50J-g_6(A)-2I.
\]
Then
\[
 T\one=2\one,\qquad T+2I\succeq0,
 \qquad 25\mathcal S_6=25(T+2I)-2J.
\]
This is the specialization of Proposition~\ref{prop:signed-complement}:
the excess parameter is \(r=42\), so the signed row sum is \(2\).
For vertices \(u,z\) at distance three, put
\[
 q_{uz}=6(A^3)_{uz}+(A^4)_{uz}.
\]
The \(2\times2\) estimate gives \(q_{uz}\in\{48,49,50,51\}\), and the kernel
argument in Theorem~\ref{thm:integral-slack} excludes \(48\).  Together with
the edge and two-path intervals in Theorem~\ref{thm:order50-feasibility}, this
gives
\[
 T_{uv}\in\{-1,0,1\}\qquad(u\ne v),\qquad T_{uu}=0.
\]
 Consequently \(T+2I\) is the Gram matrix of \(50\) vectors of norm
 \(\sqrt2\) with pairwise inner products in \(\{-1,0,1\}\) and constant signed
 row sum \(2\).  Its kernel is \(E_{-2}(A)\), so
 Proposition~\ref{prop:order50-minus-two} gives
 \[
  \operatorname{rank}(T+2I)\ge30.
 \]
 Thus the unresolved degree-six boundary is equivalently a
 root-type integral Gram problem constrained by the graph's distance classes
 and local cycle data.
\end{remark}

\begin{theorem}[Disconnected signed complement]\label{thm:order50-disconnected}
Under the hypotheses of Theorem~\ref{thm:order50-feasibility}, the underlying
signed graph of the matrix \(T\) in Remark~\ref{rem:signed-root} is
disconnected.
\end{theorem}

\begin{proof}
Let \(N_i\) denote the number of \(i\)-cycles.  In terms of the degree-six
nonbacktracking polynomials,
\[
\begin{aligned}
(g_6+2)^2={}&28144F_0+18220F_1+8838F_2+3576F_3+1233F_4\\
&+352F_5+78F_6+12F_7+F_8.
\end{aligned}
\]
Girth at least five gives \(\tr F_i(A)=0\) for \(1\le i\le4\).  For
lengths five and six, every closed nonbacktracking walk is a directed cycle.
At lengths seven and eight there are also the walks obtained by attaching a
one-edge tail to a directed five- or six-cycle.  Hence
\[
\begin{aligned}
\tr F_5(A)&=10N_5,&
\tr F_6(A)&=12N_6,\\
\tr F_7(A)&=14N_7+40N_5,&
\tr F_8(A)&=16N_8+48N_6.
\end{aligned}
\]
Removing the principal adjacency contribution and restoring the principal
signed eigenvalue \(2\) now gives
\[
\tr T^2
=8(500N_5+123N_6+21N_7+2N_8-604100).
\]
If \(P_+\) and \(P_-\) are the numbers of positive and negative signed edges,
then the signed row sum and the zero diagonal give
\[
P_+-P_-=50,\qquad
\tr T^2=2(P_++P_-)=100+4P_-.
\]
It follows that \(P_-\) is odd.

Suppose that \(T\) were connected.  By
Proposition~\ref{prop:order50-minus-two},
\(\operatorname{rank}(T+2I)\ge30\).  The root-system representation theorem for connected
edge-signed graphs with smallest eigenvalue at least \(-2\)
\cite[Theorem~2]{GreavesEtAl2015} represents \(T+2I\) in a root system of
type \(D_m\) or \(E_8\).  The rank bound excludes \(E_8\).  Thus
\[
B^{\mathsf T}B=T+2I
\]
for an integral matrix \(B\) whose columns \(b_u\) are roots
\(\pm e_i\pm e_j\).  Put \(s=B\one\).  Since
\((T+2I)\one=4\one\),
\[
b_u\mathbin{\cdot}s=4\quad\text{for every }u,
\qquad
\lVert s\rVert^2=200.
\]
Changing signs of coordinate axes if necessary, assume \(s_i\ge0\).
The coordinate-support multigraph is connected, since otherwise its columns
would split into two Gram-orthogonal families and \(T\) would be disconnected.
If \(v\) coordinates are used, then
\[
30\le v\le51:
\]
the lower bound is the Gram rank, and a connected support multigraph with
fifty root-edges has at most fifty-one vertices.

For every root, two coordinate levels satisfy
\(\pm s_i\pm s_j=4\).  Connectivity therefore places all levels in one of
\[
0,4,8,\ldots;\qquad
2,6,10,\ldots;\qquad
1,3,5,\ldots.
\]
The first family would make \(\lVert s\rVert^2\) divisible by \(16\), contrary
to \(200\equiv8\pmod {16}\).

In the second family, a level at least \(10\), together with the other
\(v-1\ge29\) levels, would contribute at least
\(10^2+29\cdot2^2>200\).  Thus only levels \(2\) and \(6\) occur.  If their
multiplicities are \(n_2,n_6\), then
\[
n_2+9n_6=50,\qquad
v=n_2+n_6=50-8n_6\ge30,
\]
so
\[
(n_2,n_6)\in\{(50,0),(41,1),(32,2)\}.
\]
Every root is now either \(e_i+e_j\) with \(s_i=s_j=2\), or
\(e_i-e_j\) with \(s_i=6\) and \(s_j=2\), after ordering the two coordinates.
Its sign pattern is unique on a fixed support, while duplicate columns would
have inner product \(2\); hence distinct roots share at most one coordinate
in this restricted family.  At a level-two coordinate let \(p_i,m_i\) be its
positive and negative incidence counts.  Then \(p_i-m_i=2\), and every
level-six coordinate has six incidences, all negative at level two.  A
negative Gram product occurs precisely when two roots meet at one level-two
coordinate with opposite incidence signs, so it is counted exactly once in
\[
P_-=\sum_{i:s_i=2}p_im_i.
\]
Consequently
\[
P_-\equiv\sum_{i:s_i=2} p_im_i
\equiv\sum_i m_i
=6n_6
\equiv0\pmod2,
\]
contradicting the parity already proved.

It remains to exclude odd levels.  Let \(A_t,B_t\) count coordinates at
levels \(4t+1,4t+3\), respectively.  Flow balance along the two level chains
gives the following identity.  Difference roots cancel within their chain
when signed incidences are summed, while the positive roots between levels
\(1\) and \(3\), the only cross-chain type, contribute equally to both sides:
\[
\sum_{t\ge0}(4t+1)A_t
=\sum_{t\ge0}(4t+3)B_t.
\]
Using this identity with \(\lVert s\rVert^2=200\) yields
\[
\frac{200-3v}{32}
=\sum_{t\ge1}\binom{t+1}{2}(A_t+B_t).
\]
Thus \(v\equiv24\pmod {32}\), which is impossible for
\(30\le v\le51\).  All three level families are excluded, proving that \(T\)
is disconnected.
\end{proof}

\section{Distance spectra of punctured Moore graphs}\label{sec:punctures}

Let \(M\) be a degree-\(k\) Moore graph of diameter two and put
\(\Delta=\sqrt{4k-3}\).

\begin{theorem}[One deleted vertex]\label{thm:one-puncture}
For \(H=M-v\),
\[
 |V(H)|=k^2,
 \qquad \delta^*(H)=k-\frac1k,
 \qquad \lambda_{\min}(D(H))=-2-\sqrt{k}.
\]
The complete distance spectrum is
\begin{align*}
\Spec D(H)={}&\{\rho_+,\rho_-\}
 \cup\{(-2+\sqrt{k})^{(k-1)},(-2-\sqrt{k})^{(k-1)}\}\\
&\cup\left\{\left(-\frac{\Delta+3}{2}\right)^{(m_+)},
\left(\frac{\Delta-3}{2}\right)^{(m_-)}\right\},
\end{align*}
where
\[
 \rho_\pm=k^2-2\pm\sqrt{k(k^3-2k^2+3k-1)},
\]
\[
 m_\pm=\frac{k(k-2)(\Delta\pm1)}{2\Delta}.
\]
Thus
\[
 \Phi(H)=k-\frac1k-2-\sqrt{k},
\]
which is positive exactly for integers \(k\ge5\).
\end{theorem}

\begin{theorem}[Endpoints of an edge]\label{thm:edge-puncture}
Let \(uv\in E(M)\) and \(H=M-\{u,v\}\).  For \(k\ge3\),
\[
 |V(H)|=k^2-1,
 \qquad \delta^*(H)=k-\frac2k,
 \qquad \lambda_{\min}(D(H))=-2-\sqrt{k}.
\]
The complete distance spectrum is
\begin{align*}
\Spec D(H)={}&\{\sigma_+,\sigma_-,k-4\}\\
&\cup\{(-2+\sqrt{k})^{(2k-4)},(-2-\sqrt{k})^{(2k-4)}\}\\
&\cup\left\{\left(-\frac{\Delta+3}{2}\right)^{(a_+)},
\left(\frac{\Delta-3}{2}\right)^{(a_-)}\right\},
\end{align*}
where
\[
 \sigma_\pm=k^2-3\pm\sqrt{k^4-2k^3+3k^2-8k+7},
\]
\[
 a_+=\frac{(k-2)(k+(k-2)\Delta)}{2\Delta},
 \quad
 a_-=\frac{(k-2)((k-2)\Delta-k)}{2\Delta}.
\]
Thus
\[
 \Phi(H)=k-\frac2k-2-\sqrt{k},
\]
which is positive exactly for integers \(k\ge5\).
\end{theorem}

\begin{proof}[Proof architecture for Theorems~\ref{thm:one-puncture} and
\ref{thm:edge-puncture}]
The displayed dual-degree formulas are the cases \(s=1\) and \(s=2\) of
Theorem~\ref{thm:small-puncture}; its proof below is independent of the
spectral decompositions.  For one deleted vertex, put \(A=N(v)\),
\(B=\Gamma_2(v)\), let \(C\) be
the \(A\)-by-\(B\) incidence matrix, and let \(B_0\) be the adjacency
matrix on \(B\).  Only pairs inside \(A\) lose their unique length-two path.
For distinct \(a,a'\in A\), choose a neighbour \(b\in B\) of \(a\).  The
vertices \(b,a'\) are nonadjacent, and their unique common neighbour \(c\)
lies in \(B\): it is not \(v\), since \(b\not\sim v\), and it is not in
\(A\), since an edge inside \(A\) would form a triangle through \(v\).
Hence \(a-b-c-a'\) is a surviving path, and therefore
\[
 D(H)=\begin{pmatrix}3(J-I)&2J-C\\2J-C^{\mathsf T}&2(J-I)-B_0\end{pmatrix}.
\]
The normalized constant quotient is
\[
 Q_1=\begin{pmatrix}
 3(k-1)&(2k-1)\sqrt{k-1}\\
 (2k-1)\sqrt{k-1}&2k^2-3k-1
 \end{pmatrix},
\]
whose eigenvalues are \(\rho_\pm\).  The identities
\(CC^{\mathsf T}=(k-1)I\), \(CB_0=J-C\), and the bottom-right block of the
Moore relation yield
\[
 \mathbb R^{V(H)}=W_{\mathrm{const}}\perp W_{\mathrm{inc}}\perp\ker C,
 \qquad 2+2(k-1)+k(k-2)=k^2.
\]
On each zero-sum incidence module the action matrix is
\(\bigl(\begin{smallmatrix}-3&-(k-1)\\-1&-1\end{smallmatrix}\bigr)\),
giving \(-2\pm\sqrt{k}\).  On \(\ker C\),
\(B_0^2+B_0-(k-1)I=0\), \(D=-2I-B_0\), and the trace of \(B_0\) is
zero after removing the constant and incidence images; dimension and trace give
\(m_\pm\).

For an adjacent deleted pair, put
\(A=N(u)\setminus\{v\}\), \(B=N(v)\setminus\{u\}\), and let \(C\) be
the residual cell.  Only pairs inside \(A\) or inside \(B\) lose a length-two path.  For
distinct \(a,a'\in A\), choose a residual neighbour \(c\) of \(a\).  The
vertices \(c,a'\) are nonadjacent, and their unique common neighbour is also
residual: it cannot be one of the deleted endpoints, cannot lie in \(A\) by
triangle-freeness, and cannot lie in \(B\) because no edge joins \(A\) to
\(B\).  This gives a surviving length-three path; the argument for \(B\) is
symmetric.  The antisymmetric constant line has eigenvalue
\(k-4\), and the normalized symmetric quotient is
\[
 Q_2=\begin{pmatrix}
 5k-8&\sqrt{(k-1)(2k-3)(4k-6)}\\
 \sqrt{(k-1)(2k-3)(4k-6)}&2k^2-5k+2
 \end{pmatrix},
\]
with eigenvalues \(\sigma_\pm\).  Two zero-sum incidence modules and a
residual kernel complete the decomposition, with
\[
 3+2(2k-4)+(k-2)^2=k^2-1.
\]
The residual adjacency trace is \(k-2\), giving \(a_\pm\).  The residual
negative root is greater than \(-2-\sqrt{k}\).  The same is true of the
smaller constant-quotient roots: in each case the leading diagonal entry of
the shifted quotient is positive, and
\[
 \det(Q_1+(2+\sqrt{k})I)
 =k(2k^2+2k^{3/2}-3k+2)>0,
\]
\[
 \det(Q_2+(2+\sqrt{k})I)
 =(\sqrt{k}-1)(\sqrt{k}+1)
 (2k^2+2k^{3/2}-3k+2\sqrt{k}+6)>0.
\]
Exact quotient normalization, trace-to-multiplicity equations, replacement
paths, and all least-root comparisons are independently checked by
\codefile{scripts/verify_proof_audit_05_small_moore_punctures.py}.
\end{proof}

\begin{theorem}[Two nonadjacent deleted vertices]\label{thm:nonadjacent-puncture}
Let \(u,v\) be nonadjacent vertices of \(M\), \(k\ge5\), and put
\(H=M-\{u,v\}\).  Define
\begin{align*}
R_k(x)={}&x^4+(10-2k^2)x^3+(2k^3-17k^2-2k+36)x^2\\
&+(12k^3-49k^2-4k+53)x\\
&-2k^4+17k^3-38k^2+5k+20,
\end{align*}
\[
 M_- =\frac{k(k-2)+(k^2-4k+2)\Delta}{2\Delta},
 \quad
 M_+ =\frac{-k(k-2)+(k^2-4k+2)\Delta}{2\Delta}.
\]
Then
\begin{align*}
\chi_{D(H)}(x)={}&(x-k+3)R_k(x)
 (x^2+4x-k+3)^{k-2}\\
&\cdot(x^2+4x-k+5)^{k-2}
 \left(x+\frac{\Delta+3}{2}\right)^{M_-}\\
&\cdot\left(x-\frac{\Delta-3}{2}\right)^{M_+}.
\end{align*}
Moreover,
\[
 \delta^*(H)=k-\frac2k,
\]
and \(H\) is a strict counterexample for every realizable integer \(k\ge6\).
\end{theorem}

\begin{proof}
The deleted vertices have a unique common neighbour \(w\).  With
\(A=N(u)\setminus\{w\}\), \(B=N(v)\setminus\{w\}\),
\(C=N(w)\setminus\{u,v\}\), and residual cell \(Z\), the surviving graph
has an equitable five-cell geometry and a perfect matching between \(A\) and
\(B\).  Explicit length-three paths show that every pair whose unique
length-two path was deleted has new distance exactly three.  The five-cell
distance quotient has characteristic polynomial \((x-k+3)R_k(x)\).

Let \(R_A,R_B,R_C\) be the incidence matrices from the three nonconstant
cells to \(Z\), and let \(T\) be the adjacency matrix on \(Z\).  The Moore
identity splits the remaining space into matched symmetric and antisymmetric
modules, a common-neighbour module, and
\(K=\ker R_A\cap\ker R_B\cap\ker R_C\).  The first three modules give
the two displayed quadratic factors and \(k-3\) copies of each Moore linear
factor.  On \(K\),
\[
 T^2+T-(k-1)I=0,
 \qquad D=-2I-T.
\]
The constant direction of \(Z\) has \(T\)-eigenvalue \(k-3\), the
symmetric, antisymmetric, and common-neighbour images have eigenvalues
\(-2,0,-1\) with dimensions \(k-2,k-2,k-3\).  Since \(T\) has zero
diagonal, its total trace is zero, and therefore
\(\tr(T|_K)=2(k-2)\).  Dimension and trace now give the residual
multiplicities; after adding the common-neighbour copies they are precisely
\(M_-\) and \(M_+\).  The full dimension count is
\[
 5+4(k-2)+2(k-3)+(k-2)(k-4)=k^2-1.
\]

The value \(\delta^*(H)=k-2/k\) is again the \(s=2\) case of
Theorem~\ref{thm:small-puncture}.  For strictness, the distance-increase
matrix is the adjacency matrix of two copies of \(K_k\) meeting in \(w\),
and hence has least eigenvalue
\[
 \lambda_{\min}(E_S)=\frac{k-2-\sqrt{k^2+4k-4}}2.
\]
Since the parent Moore graph has score
\(k-(3+\Delta)/2\) and deletion lowers the minimum dual degree by
\(2/k\),
Proposition~\ref{prop:deletion-stability} gives
\[
 \Phi(H)\ge
 \frac{3k-5-\Delta-\sqrt{k^2+4k-4}}2-\frac2k.
\]
For \(k\ge6\), the sign-preserving squarings reduce positivity of this
lower bound to \(P(k)>0\), where
\[
 P(k)=4k^8-39k^7+95k^6+9k^5-173k^4-36k^3+116k^2+80k+16.
\]
Writing \(k=m+6\) gives
\[
\begin{aligned}
 P(m+6)={}&4m^8+153m^7+2489m^6+22329m^5\\
 &+119437m^4+382236m^3+685268m^2\\
 &+559616m+75952>0.
\end{aligned}
\]
The direct-sum decomposition, injectivity, orthogonality, all recomputed
distances, and the polynomial sign
are independently checked by
\codefile{scripts/verify_proof_audit_03_nonadjacent_puncture.py} and
\codefile{scripts/verify_research_extensions_exact.py}.
\end{proof}

\begin{proposition}[Deletion stability]\label{prop:deletion-stability}
Let \(G\) be connected, \(S\subseteq V(G)\), and suppose \(H=G-S\) is
connected.  Put
\[
 D_0=D(G)[V(H)],
 \qquad E_S=D(H)-D_0,
\]
\[
 a=\delta^*(G),
 \quad b=\delta^*(H),
 \quad \gamma=\Phi(G).
\]
Then
\[
 \resultbox{
 \Phi(H)\ge\gamma-(a-b)+\lambda_{\min}(E_S).
 }
\]
\end{proposition}

\begin{proof}
One has
\[
 D(H)+bI=(D_0+aI)+E_S-(a-b)I.
\]
Principal-submatrix interlacing gives
\(\lambda_{\min}(D_0+aI)\ge\gamma\), and Weyl's inequality gives the result.
For Moore punctures, \(E_S\) is an explicit distance-increase graph; the exact
specializations are checked by
\codefile{scripts/verify_research_extensions_exact.py}.
\end{proof}

\section{Small punctures and exact Hoffman--Singleton robustness}
\label{sec:robustness}

\begin{theorem}[Small-puncture normal form]\label{thm:small-puncture}
Let \(M\) be a degree-\(k\) Moore graph, let \(S\subseteq V(M)\) have size
\(s\le k-1\), and put \(H=M-S\).  Then \(H\) is connected, has diameter at
most three, and
\[
 \resultbox{\delta^*(H)=k-\frac{s}{k}.}
\]
Let \(B\) be the surviving-vertex by deleted-vertex incidence matrix, with
\(B_{xz}=1\) exactly when \(x\sim_M z\).  Then
\[
 \resultbox{
 D(H)=2(J-I)-A(H)+BB^{\mathsf T}
 -\operatorname{diag}(BB^{\mathsf T}).
 }
\]
\end{theorem}

\begin{proof}
The case \(s=0\) is immediate, so assume \(1\le s\le k-1\).  Let
\(x,y\in V(H)\) be nonadjacent in \(M\), and suppose their unique common
neighbour \(z\) lies in \(S\).  For every
\(a\in N_M(x)\setminus\{z\}\), the vertices \(a,y\) are nonadjacent and
have a unique common neighbour \(b_a\).  Thus
\[
 x-a-b_a-y
\]
is a length-three path.  These \(k-1\) paths are internally vertex-disjoint:
a shared vertex \(b_a=b_{a'}\) would give a four-cycle, while
\(b_a=a'\) would give a triangle through \(x\).  Besides \(z\), at most
\(s-1\le k-2\) vertices are deleted, so one path survives.  Hence a destroyed
length-two path becomes distance exactly three, which proves the matrix formula.

For \(x\in V(H)\), let \(t_x=|N_M(x)\cap S|\).  A deleted neighbour of
\(x\) is adjacent to no surviving neighbour of \(x\), by triangle-freeness;
each other deleted vertex has at most one common neighbour with \(x\).  Thus
\[
 \sum_{y\in N_H(x)}t_y\le s-t_x.
\]
Consequently,
\[
 d_H^*(x)
 \ge k-\frac{s-t_x}{k-t_x}
 \ge k-\frac{s}{k},
\]
where the last inequality follows from
\[
 \left(k-\frac{s-t}{k-t}\right)-\left(k-\frac{s}{k}\right)
 =\frac{t(k-s)}{k(k-t)}\ge0.
\]
The intersection bound
\[
 \left|\bigcap_{z\in S}\Gamma_2(z)\right|
 \ge k^2+1-s(k+1)\ge2
\]
provides a surviving vertex \(x\) at distance two from every deleted vertex.
For \(z\in S\), let \(y_z\) be the unique common neighbour of \(x,z\).  If
\(y_z\in S\), then the choice of \(x\) would require
\(d_M(x,y_z)=2\), contradicting \(x\sim y_z\).  Thus all witnesses survive,
each deleted vertex contributes exactly once to the neighbour-degree deficit,
and equality is attained.  The
replacement paths, boundary case \(s=k-1\), distance formula, and attainment
step are independently checked by
\codefile{scripts/verify_proof_audit_12_small_puncture.py}.
\end{proof}

\begin{corollary}[Uniform deletion stability]
\label{cor:uniform-deletion}
Under the hypotheses of Theorem~\ref{thm:small-puncture}, put
\[
 t_x=|N_M(x)\cap S|,
 \qquad
 \tau(S)=\max_{x\in V(M-S)}t_x.
\]
Then
\[
 \resultbox{
 \Phi(M-S)\ge
 k-\frac{3+\sqrt{4k-3}}2-\frac{s}{k}-\tau(S).
 }
\]
In particular, the deletion is a strict counterexample whenever the
right-hand side is positive.  Since \(\tau(S)\le s\), every deletion of
\(s\) vertices is strict whenever
\[
 \resultbox{
 s\left(1+\frac1k\right)
 <k-\frac{3+\sqrt{4k-3}}2.
 }
\]
Equivalently, every deletion of at most
\[
 r_k=\min\left\{
 k-1,\
 \left\lceil
 \frac{k}{k+1}\left(k-\frac{3+\sqrt{4k-3}}2\right)
 \right\rceil-1
 \right\}
\]
vertices is strict.  In particular, \(r_7=2\); for a hypothetical
degree-\(57\) Moore graph, \(r_{57}=47\).
\end{corollary}

\begin{proof}
By Theorem~\ref{thm:small-puncture}, the distance-increase matrix is
\[
 E_S=BB^{\mathsf T}-\operatorname{diag}(t_x).
\]
Since \(BB^{\mathsf T}\succeq0\) and
\(\operatorname{diag}(t_x)\preceq\tau(S)I\), one has
\(\lambda_{\min}(E_S)\ge-\tau(S)\).  The parent Moore graph has score
\[
 \Phi(M)=k-\frac{3+\sqrt{4k-3}}2,
\]
and Theorem~\ref{thm:small-puncture} shows that deletion lowers the minimum
dual degree by \(s/k\).  Proposition~\ref{prop:deletion-stability} now gives
the first bound.  The second follows from \(\tau(S)\le s\), and solving its
strict scalar inequality for the largest integral \(s\) gives \(r_k\).
\end{proof}

\begin{theorem}[Hoffman--Singleton robustness radius]\label{thm:hs-radius}
Let \(M\) be the Hoffman--Singleton graph.  Every induced graph \(M-S\) with
\(|S|\le5\) is a strict counterexample to WOW-284.  This is sharp in the
universal sense: there exists a six-vertex set whose deletion is not strict.
Hence the universal vertex-deletion robustness radius is exactly five.
\end{theorem}

\begin{proof}
Theorem~\ref{thm:small-puncture} gives
\[
 \delta^*(M-S)=\frac{49-|S|}{7}.
\]
Corollary~\ref{cor:uniform-deletion} already proves strictness uniformly for
\(|S|\le2\).  To obtain the sharp universal radius, two explicitly stored
permutations are verified edge-by-edge as automorphisms of the coordinate
graph; the group they generate has order \(252000\).  Their orbits on
deletion sets of sizes \(0,1,2,3,4,5\) have counts
\[
 1,1,2,4,11,33.
\]
For every one of the \(52\) representatives, exact fraction arithmetic gives
\[
 7D(M-S)+(49-|S|)I\succ0.
\]
Orbit-size sums equal \(\binom{50}{s}\), so every labelled set is covered.

For sharpness, delete
\[
 \{P_{2,4},P_{3,1},P_{3,4},Q_{2,1},Q_{3,4},Q_{4,4}\}.
\]
The resulting graph has \(\delta^*=43/7\), while an exact
\(LDL^{\mathsf T}\) decomposition of \(7D+43I\) has exactly one negative and
no zero pivot.  Thus it is not strict.  Generator action, orbit exhaustion,
BFS distances, the small-puncture formula, and a handwritten rational
\(LDL^{\mathsf T}\) implementation are independently checked by
\codefile{scripts/verify_proof_audit_13_hs_robustness.py}.
\end{proof}

\begin{remark}
The theorem asserts that every deletion through size five succeeds and that at
least one deletion of size six fails.  It does not assert that every
six-vertex deletion fails.
\end{remark}

\section{Equality and obstructions to natural construction families}
\label{sec:controls}

\begin{theorem}[Equality boundary]\label{thm:equality-boundary}
Let \(G\) be connected, \(k\)-regular, of girth at least five and diameter
three.  Then
\[
 \Phi(G)=0
 \quad\Longleftrightarrow\quad
 \max_{\theta\ne k}|\theta+1|=\sqrt{2k-2}.
\]
Equivalently, \(D+kI\) is positive semidefinite and singular.  If \(2k-2\) is
not a square, the two boundary adjacency eigenvalues occur with equal
multiplicity, so the distance eigenvalue \(-k\) has even multiplicity.  If
\(2k-2\) is a square, then \(k=2r^2+1\) for some integer \(r\).
\end{theorem}

\begin{proof}
This is immediate from Theorem~\ref{thm:diameter-three-score}.  In the
nonsquare case, the two boundary values are algebraic conjugates, so their
multiplicities in the integral characteristic polynomial agree.  In the square
case, \(2k-2=s^2\) forces \(s\) even.  The exact scalar audit is
\codefile{scripts/verify_equality_boundary.py}.
\end{proof}

J{\o}rgensen's \(9\)-regular order-\(96\) graph of girth five is an exact equality
case:
\[
 \delta^*=9,
 \qquad \lambda_{\min}(D)=-9,
\]
with multiplicity eight.  The construction is due to J{\o}rgensen
\cite{Jorgensen2005}; the three local graph representations, handwritten
graph6 decoder, characteristic polynomials, root intervals, and provenance
boundary are checked by
\codefile{scripts/verify_proof_audit_09_jorgensen96.py}.

\subsection{A prime-field obstruction}

For an odd prime \(q\ge7\) and \(1\le m\le q\), define \(G(q,m)\) on
vertices \(P_{i,j},Q_{k,\ell}\), with \(0\le i,k<m\) and
\(j,\ell\in\mathbb F_q\), by
\begin{align*}
P_{i,j}&\sim P_{i,j\pm1},\\
Q_{k,\ell}&\sim Q_{k,\ell\pm2},\\
P_{i,j}&\sim Q_{k,ik+j}.
\end{align*}
This is a balanced specialization of known finite-field girth-five
constructions \cite{AbreuEtAl2008}.  It is \((m+2)\)-regular.  A coordinate
common-neighbour calculation shows that it has no triangle or \(4\)-cycle for
\(q\ge7\): for a cross pair, the possible same-side common neighbours require
residues in the disjoint sets \(\{\pm1\}\) and \(\{\pm2\}\).

\begin{theorem}\label{thm:prime-field}
If \(G(q,m)\) has diameter three, then it is not a strict counterexample to
WOW-284.
\end{theorem}

\begin{proof}
The zero Fourier block has eigenvalues \(m+2\), \(2-m\), and \(2\) with
multiplicity \(2m-2\); the shifted WOW window leaves only
\(m\in\{4,5,6\}\).  Let \(\omega=e^{2\pi\mathrm i/q}\).  On the nonzero
character \(t=1\), the adjacency block has form
\[
 \begin{pmatrix}aI&M\\M^*&bI\end{pmatrix},
 \qquad a=2\cos\frac{2\pi}{q},
 \quad b=2\cos\frac{4\pi}{q},
 \quad M_{ik}=\omega^{ik}.
\]
If \(\sigma\) is a singular value of \(M\), the associated two-dimensional
invariant subspace carries
\(\bigl(\begin{smallmatrix}a&\sigma\\\sigma&b\end{smallmatrix}\bigr)\).
Since \(\|M\|_F^2=m^2\), one singular value satisfies
\(\sigma^2\ge m\), and the block has a nonprincipal eigenvalue at least
\[
 \sqrt m+\frac{a+b}{2}
 \ge \sqrt m+\cos(\pi/7)-\frac12.
\]
Here the last inequality uses
\(2\cos(2\pi/7)+2\cos(4\pi/7)=2\cos(\pi/7)-1\).
Moreover \(\cos(\pi/7)>\sqrt3/2\), and the increasing function
\(h(m)=\sqrt{2m+2}-\sqrt m\) satisfies
\[
 h(m)\le h(6)=\sqrt{14}-\sqrt6
 <\frac{27}{20}<\frac{\sqrt3+1}{2}.
\]
Thus the displayed nonprincipal eigenvalue lies above
\(-1+\sqrt{2m+2}\) for \(m=4,5,6\).  The common-neighbour
case split, Fourier reduction, radical comparisons, and exact \(q=7\) controls
are independently checked by
\codefile{scripts/verify_proof_audit_08_prime_field.py}.
\end{proof}

\subsection{Layer-respecting matching deletions}

For \(\pi\in S_5\), delete the perfect matching
\[
 \mathcal M_\pi=\bigl\{\{P_{i,j},Q_{\pi(i),i\pi(i)+j}\}:i,j\in\F\bigr\}
\]
from the Hoffman--Singleton graph.

\begin{theorem}\label{thm:matching-deletions}
Each deletion produces a connected simple \(6\)-regular graph of order
\(50\), girth five, and diameter four.  The \(120\) labelled graphs form
exactly two isomorphism classes.  The \(20\) affine permutations have
\[
 \lambda_{\min}(D)=-13,
 \qquad \Phi=-7,
\]
and the \(100\) nonaffine permutations have
\[
 \lambda_{\min}(D)=-6-\sqrt{61},
 \qquad \Phi=-\sqrt{61}.
\]
Thus every member strictly satisfies the conjectured inequality.
\end{theorem}

\begin{proof}
Every \(P_{i,j}\) occurs once in \(\mathcal M_\pi\); for a fixed
\(Q_{k,\ell}\), the unique incident matching edge is obtained from
\(i=\pi^{-1}(k)\) and \(j=\ell-ik\).  Thus \(\mathcal M_\pi\) is a
perfect matching, and its deletion leaves a simple \(6\)-regular graph.
Deleting edges cannot create a short cycle, while the same-layer pentagons
remain; exact breadth-first search gives connectedness and diameter four.

Explicit type-preserving and type-swapping coordinate automorphisms generate
orbits of sizes \(20\) and \(100\).  The representatives have different
adjacency characteristic polynomials, so the orbits are distinct isomorphism
classes.  Exact distance characteristic polynomials and Sturm separators give
the stated least roots.  All \(120\) matchings, \(400\) coordinate maps,
\(48{,}000\) matching images, orbit coverage, graph hypotheses, and root
certificates are checked
by \codefile{scripts/verify_proof_audit_07_layer_matchings.py}.
\end{proof}

\section{Exact computation and formal verification}\label{sec:verification}

The analytic arguments above are primary.  Exact computation is used in three
roles: to certify explicitly labelled finite graphs, to exhaust precisely
specified finite orbit families, and to check symbolic identities whose
written derivations are also supplied.  No theorem-level sign or eigenvalue
ordering uses floating-point arithmetic.  The exact symbolic and graph
computations use SymPy and NetworkX \cite{MeurerEtAl2017,HagbergSchultSwart2008}.

The accompanying release archives the labelled graph data, exact rational
and polynomial certificates, and exhaustive finite-orbit records used here.
These materials make the finite computations independently reproducible
without interrupting the mathematical exposition with a file-by-file index.

The explicit \(50\)-vertex counterexample is fully formalized in Lean 4.31
with Mathlib 4.31 \cite{deMouraUllrich2021,Mathlib2020}.  The development
checks the coordinate graph, regularity, the exhaustive common-neighbour
certificate, girth five, the adjacency-square and distance-matrix identities,
and an exact rational diagonalization with multiplicities.  Thus the
\(50\)-vertex result is verified at graph level, including its least distance
eigenvalue and strict WOW-284 gap.

Lean also kernel-checks finite spectral certificates attached to the explicit
constructions of orders \(38,39,40,42\).  For orders \(38,39,42\), the
endpoints certify the stated minimum dual-degree data, positive definiteness of
the corresponding shifted finite distance matrix, and hence the strict bound
for every nonzero real eigenpair.  For order \(40\), Lean certifies a two-sided
invertible exact diagonalization, the diagonal-entry multiplicities, the
attained least eigenvalue \(-5\), dual degree six, and gap one.  These
non-\(50\) endpoints are finite spectral certificates: they do not identify
each semantic matrix with \texttt{SimpleGraph.dist} and bundle every
structural hypothesis into one public theorem.

Separately, Lean formalizes the analytic optimization statement in
Theorem~\ref{thm:lp-ceiling}.  For every integer \(k\ge4\) and every admissible
finitely supported expansion
\[
  f=\sum_i c_iF_i,
  \qquad c_0>0,
  \qquad c_i\ge0\quad(i\ge5),
  \qquad f\vert_{I_k}\le0,
\]
it proves
\[
  B_kc_0\le f(k),
  \qquad B_k=\frac{(k+2)(k^2+3)}6.
\]
The formal development defines the coefficient family of the displayed
quartic \(f_*\), proves that it is admissible and attains equality, and proves
that every equality case is its unique positive scalar multiple, both as a
polynomial and at coefficient level.  The kernel-checked dependency chain
includes the nonbacktracking recurrence, the exact primal expansion, positivity
and moment identities for the three-point dual certificate, strict slacks in
degrees \(5\) through \(9\), the Chebyshev tail for every degree at least
\(10\), finite-support weak duality, complementary slackness, and optimizer
rigidity.

This LP formalization is deliberately graph-independent.  It does not
formalize the trace interpretation of the \(F_i(A)\), the girth-five
vanishing and nonnegativity statements, the spectral trace decomposition, or
the passage from the analytic LP inequality to graph-order bounds.
Consequently the graph conclusions \(n<B_k\), the degree-six reductions
\(n\le51\) and \(n\le50\), the integral optimal-slack theorem, the
 three-to-one excess bound and equality rigidity, the sharper low-degree
 windows, the order-\(50\) signed-complement disconnection, and the local Gram
 hierarchy remain conventional analytic proofs, with exact executable audits
 where cited.  The punctured-Moore
spectra, the general deletion-stability inequality, the small-puncture normal
form, the uniform Gram bound, and the Hoffman--Singleton deletion-radius
theorem are likewise analytic results supported by exact Python audits; they
are not part of the Lean claim.

\medskip
\noindent
The public Lean development is sorry-free and kernel-checked by Lean 4.31
within the scope stated above.

\section{Broader mathematical questions}\label{sec:scope}

The mechanisms in this paper point beyond the conjecture that motivated them.
They connect extremal polynomial certificates, local positive-semidefinite
constraints, perturbations of metric operators, exact computation, and formal
proof.  At that methodological level, they suggest the following questions.

\begin{enumerate}
\item Given an optimal polynomial or semidefinite certificate for a global
spectral inequality, when do the principal minors of its slack operator form a
complete hierarchy of local constraints?  Can rank or flat-extension
conditions force finite convergence and reconstruct the extremal discrete
objects?
\item Is there a general inverse theory for spectral inequalities in which a
small global defect forces an object to lie quantitatively near a structured
algebraic model?  What notions of distance and stability are natural for such
statements?
\item Which one-variable spectral arguments survive when homogeneity or
commutativity is lost?  Can degree-weighted operators, nonbacktracking
operators, or matrix-valued orthogonal polynomials provide comparably sharp
certificates for irregular or multitype systems?
\item Can a unified perturbation-and-limit theory for globally defined metric
operators explain when local modifications preserve a spectral inequality
and distinguish sporadic finite phenomena from finite shadows of infinite
families?
\end{enumerate}

OpenAI ChatGPT-5.6 Sol Pro assisted with adversarial proof checking, proof
exploration, and Lean formalization.  The author assumes full responsibility
for the mathematics, attribution, and conclusions.  The source, exact
certificates, and build instructions are available at
\href{https://github.com/SamPetkov/wow284}{\texttt{github.com/SamPetkov/wow284}}
and correspond to release \texttt{v2.3.0}.

\clearpage
\begingroup
\footnotesize
\setlength{\bibsep}{0pt}
\renewcommand{\bibsection}{\section*{References}}
\bibliographystyle{unsrtnat}
\bibliography{references}
\endgroup

\end{document}
